% =====================================================================
%  numerical-scheme-ru.tex
%
%  Русская версия numerical-scheme.tex.  Структура, метки, формулы и
%  имена в коде совпадают один в один, поэтому \ref/\eqref переносятся
%  между двумя файлами без правок.
%
%  Движок BF_v_no_resampling_v2.py: базовый алгоритм, десять трюков,
%  из которых он собран, и измерительный слой, превращающий прогон в число.
%
%  ДВА СПОСОБА ИСПОЛЬЗОВАНИЯ
%  -------------------------
%  (a) Отдельно: компилируется как есть.  Файл самодостаточен и не
%      ссылается ни на какую внешнюю библиографию, так что undefined
%      references не будет.
%  (b) Как Part внутри main.tex: удалить всё между строками
%      "ПРЕАМБУЛА -- ОТРЕЗАТЬ ОТСЮДА" и "ПРЕАМБУЛА -- ОТРЕЗАТЬ ДОСЮДА",
%      удалить финальный \end{document}, положить в Parts/ и добавить
%          % =====================================================================
%  numerical-scheme-ru.tex
%
%  Русская версия numerical-scheme.tex.  Структура, метки, формулы и
%  имена в коде совпадают один в один, поэтому \ref/\eqref переносятся
%  между двумя файлами без правок.
%
%  Движок BF_v_no_resampling_v2.py: базовый алгоритм, десять трюков,
%  из которых он собран, и измерительный слой, превращающий прогон в число.
%
%  ДВА СПОСОБА ИСПОЛЬЗОВАНИЯ
%  -------------------------
%  (a) Отдельно: компилируется как есть.  Файл самодостаточен и не
%      ссылается ни на какую внешнюю библиографию, так что undefined
%      references не будет.
%  (b) Как Part внутри main.tex: удалить всё между строками
%      "ПРЕАМБУЛА -- ОТРЕЗАТЬ ОТСЮДА" и "ПРЕАМБУЛА -- ОТРЕЗАТЬ ДОСЮДА",
%      удалить финальный \end{document}, положить в Parts/ и добавить
%          % =====================================================================
%  numerical-scheme-ru.tex
%
%  Русская версия numerical-scheme.tex.  Структура, метки, формулы и
%  имена в коде совпадают один в один, поэтому \ref/\eqref переносятся
%  между двумя файлами без правок.
%
%  Движок BF_v_no_resampling_v2.py: базовый алгоритм, десять трюков,
%  из которых он собран, и измерительный слой, превращающий прогон в число.
%
%  ДВА СПОСОБА ИСПОЛЬЗОВАНИЯ
%  -------------------------
%  (a) Отдельно: компилируется как есть.  Файл самодостаточен и не
%      ссылается ни на какую внешнюю библиографию, так что undefined
%      references не будет.
%  (b) Как Part внутри main.tex: удалить всё между строками
%      "ПРЕАМБУЛА -- ОТРЕЗАТЬ ОТСЮДА" и "ПРЕАМБУЛА -- ОТРЕЗАТЬ ДОСЮДА",
%      удалить финальный \end{document}, положить в Parts/ и добавить
%          % =====================================================================
%  numerical-scheme-ru.tex
%
%  Русская версия numerical-scheme.tex.  Структура, метки, формулы и
%  имена в коде совпадают один в один, поэтому \ref/\eqref переносятся
%  между двумя файлами без правок.
%
%  Движок BF_v_no_resampling_v2.py: базовый алгоритм, десять трюков,
%  из которых он собран, и измерительный слой, превращающий прогон в число.
%
%  ДВА СПОСОБА ИСПОЛЬЗОВАНИЯ
%  -------------------------
%  (a) Отдельно: компилируется как есть.  Файл самодостаточен и не
%      ссылается ни на какую внешнюю библиографию, так что undefined
%      references не будет.
%  (b) Как Part внутри main.tex: удалить всё между строками
%      "ПРЕАМБУЛА -- ОТРЕЗАТЬ ОТСЮДА" и "ПРЕАМБУЛА -- ОТРЕЗАТЬ ДОСЮДА",
%      удалить финальный \end{document}, положить в Parts/ и добавить
%          \include{Parts/numerical-scheme-ru}
%      рядом с уже имеющимся \include{Parts/unified-cascade_v4}.
%      В преамбулу main.tex тогда нужно добавить:
%          \usepackage{algorithm}
%          \usepackage{algpseudocode}
%          \usepackage{booktabs}
%      и, поскольку main.tex сейчас английский, переключить babel на
%          \usepackage[english,russian]{babel}
%      и кодировку на \usepackage[T1,T2A]{fontenc}.
% =====================================================================

% ------------------------------------------ ПРЕАМБУЛА -- ОТРЕЗАТЬ ОТСЮДА
\documentclass[a4paper,12pt]{article}

\usepackage[T1,T2A]{fontenc}
\usepackage[utf8]{inputenc}
\usepackage[english,russian]{babel}
\usepackage{amsmath,amssymb}
\usepackage{geometry}
\geometry{margin=2.5cm}
\usepackage{graphicx}
\usepackage{booktabs}
\usepackage{algorithm}
\usepackage{algpseudocode}
\usepackage[unicode,colorlinks=true,linkcolor=blue,citecolor=blue,urlcolor=blue]{hyperref}

\newcommand{\minj}{m_{\mathrm{inj}}}
\newcommand{\msink}{m_{\mathrm{sink}}}
\newcommand{\fmaj}{f^{\mathrm{maj}}}
\newcommand{\code}[1]{\texttt{#1}}
\emergencystretch=2em   % русские слова длиннее английских

\begin{document}
\selectlanguage{russian}

\title{Численная схема: метод Монте-Карло с мажорантной частотой\\ для каскада в пространстве масс}
\author{}
\date{}
\maketitle
% ------------------------------------------ ПРЕАМБУЛА -- ОТРЕЗАТЬ ДОСЮДА

\section{Область применимости и принцип построения}
\label{sec:num_scope}

Движок интегрирует стохастический каскад в пространстве масс в четырёх
конфигурациях,
\[
  \text{процесс} \in \{\text{коагуляция},\ \text{фрагментация}\},
  \qquad
  \text{система} \in \{\text{замкнутая},\ \text{открытая}\},
\]
одним и тем же путём в коде. Между четырьмя случаями различаются ровно две
вещи:

\begin{enumerate}
  \item что принятое событие делает с двумя выбранными частицами, и
  \item граничные условия --- инжекция плюс поглощающий сток, либо ни того,
        ни другого.
\end{enumerate}

Всё остальное --- построение мажоранты, выбор пары, часы, учёт ---
общее. Это сознательное ограничение, а не удобство: расхождение между
двумя из четырёх случаев тогда нельзя списать на два разных интегратора.

Описываемый здесь вариант не несёт весов частиц. Одна расчётная частица
есть одна физическая частица, $w \equiv 1$ на весь прогон, поэтому
$\sum_i w_i$ и $\sum_i w_i m_i$ сохраняются с машинной точностью на каждом
шаге. Измеряемый дрейф массы тождественно равен нулю, и любое ненулевое
значение --- это ошибка в коде, а не статистическое отклонение.
Раздел~\ref{sec:num_cost} говорит, чего это стоит.

\subsection{Состояние}

Состояние --- четыре параллельных массива по живым частицам плюс разбиение
пространства масс на бины:

\begin{center}
\small
\begin{tabular}{@{}ll@{}}
\toprule
\code{mass[i]}       & масса частицы $m_i$ \\
\code{inj\_time[i]}  & время рождения, нужно для изохрон (разд.~\ref{sec:trick7}) \\
\code{bin\_of[i]}    & номер массового бина, в котором лежит $i$ \\
\code{pos\_in\_bin[i]} & позиция $i$ внутри списка индексов этого бина \\
\midrule
\code{edges[0..B]}   & $B+1$ логарифмически расположенных границ бинов \\
\code{bins[b]}       & список индексов частиц в бине $b$ \\
\code{N\_bin[b]}     & заселённость $N_b$ \\
\bottomrule
\end{tabular}
\end{center}

Частица массы $m$ попадает в бин
$b(m) = \max\bigl(0,\min(B-1,\ \mathrm{searchsorted}(\text{edges},m)-1)\bigr)$.
Зажим (clamping) в этом выражении --- ловушка, и она разобрана в
разд.~\ref{sec:num_pitfalls}.

\section{Базовый алгоритм}
\label{sec:num_algorithm}

\subsection{Инициализация}

\begin{algorithm}[h]
\caption{Инициализация: бины, таблица мажоранты, вектор скоростей}
\label{alg:setup}
\begin{algorithmic}[1]
\State разложить все начальные частицы по бинам; сформировать заселённости $N_b$
\State \textbf{таблица мажоранты:} для всех $b,c$ \quad
       $\fmaj_{bc} \gets K(\text{edges}[b+1],\ \text{edges}[c+1])$
       \Comment{верхние углы}
\State $T_b \gets \sum_c \fmaj_{bc} N_c$
\State $\text{row}_b \gets N_b\,(T_b - \fmaj_{bb})$
\State $R \gets \sum_b \text{row}_b$
\end{algorithmic}
\end{algorithm}

Угловое правило $\fmaj_{bc}=K(\text{edges}[b+1],\text{edges}[c+1])$
ограничивает $K$ сверху на всей паре бинов \emph{при условии, что $K$ не
убывает по обоим аргументам}. Это единственное структурное требование к
ядру, и оно проверяется во время счёта, а не предполагается: каждая
принятая попытка проверяет $K(m_i,m_j) \le \fmaj_{bc}$ и падает, если это
нарушено.

По построению
\begin{equation}
  R \;=\; \sum_b N_b\bigl(T_b - \fmaj_{bb}\bigr)
    \;=\; \sum_{i \neq j} \fmaj_{ij}
    \;=\; 2 \sum_{i<j} \fmaj_{ij},
  \label{eq:num_R}
\end{equation}
то есть $R$ пробегает \emph{упорядоченные} пары. Двойка в
уравнении~\eqref{eq:num_R} --- источник двойки в часах,
уравнение~\eqref{eq:num_clock}.

\subsection{Главный цикл}

\begin{algorithm}[h]
\caption{Главный цикл: одна мажорантная попытка}
\label{alg:mainloop}
\begin{algorithmic}[1]
\While{ни одно условие остановки не выполнено}
  \State $\Delta t \gets 2V/(wR)$; \quad $t \gets t + \Delta t$
         \Comment{часы, разд.~\ref{sec:trick5}}
  \State \textbf{инжекция (открытые системы):}
         $\nu \gets \nu + q\,\Delta t / w$; \quad
         $k \gets \lfloor \nu \rfloor$; \quad $\nu \gets \nu - k$;
         добавить $k$ частиц массы $\minj$
  \State выбрать бин-строку $b$ с $P(b) \propto \text{row}_b$
  \State выбрать бин-столбец $c$ с
         $P(c) \propto N_c \fmaj_{bc}$, беря $(N_b-1)\fmaj_{bb}$ при $c=b$
  \State выбрать $i$ равномерно из \code{bins[b]}, $j \neq i$ равномерно из \code{bins[c]}
  \State \textbf{прореживание:} принять с вероятностью $K(m_i,m_j)/\fmaj_{bc}$;
         иначе \textbf{continue} \Comment{нулевое событие}
  \State применить событие (алг.~\ref{alg:event}); применить поглощающий сток
  \State записать снапшот, если этого требует расписание (разд.~\ref{sec:trick10})
\EndWhile
\end{algorithmic}
\end{algorithm}

Отклонённая попытка --- не ошибка: именно прореживание делает доминирующий
процесс точным. Доля принятия, выводимая как \code{acc} в логе прогона,
есть $O(1)$ по построению бин-парной таблицы --- измерено $0{,}95$ для
геометрического ядра на бинах шириной $0{,}1$~декады, где угловой перебор
составляет $10^{0.1\lambda}$ на аргумент.

\subsection{Событие}

\begin{algorithm}[h]
\caption{Принятое событие --- единственное место, где четыре случая расходятся}
\label{alg:event}
\begin{algorithmic}[1]
\If{коагуляция}
  \State $m_{\mathrm{new}} \gets m_i + m_j$
  \State унаследовать возраст $i$ от пары $(i,j)$ по \code{age\_rule} (разд.~\ref{sec:trick7})
  \State $m_i \gets m_{\mathrm{new}}$; перебиновать $i$; \textbf{удалить} $j$
  \State \textbf{если} открытая \textbf{и} $m_{\mathrm{new}} \ge \msink$ \textbf{то}
         удалить $i$; $M_{\mathrm{out}} \mathrel{+}= w\,m_{\mathrm{new}}$;
         $n_{\mathrm{out}} \mathrel{+}= 1$
\Else \Comment{фрагментация}
  \State $p \gets \arg\max(m_i,m_j)$; \quad $s \gets \arg\min(m_i,m_j)$
  \State \textbf{если} $m_s < f\,m_p$ \textbf{то continue}
         \Comment{порог разрушения, разд.~\ref{sec:trick8}}
  \State $\xi \sim U(0,1)$; \quad $m_a \gets \xi m_p$, \ $m_b \gets (1-\xi)m_p$
  \State оба осколка наследуют $t_{\mathrm{born}}$ единственного родителя $p$
  \State $m_p \gets m_a$; перебиновать $p$; \textbf{добавить} частицу массы $m_b$
  \State \textbf{если} открытая: удалить каждый осколок с $m \le \msink$,
         начиная со старшего индекса
\EndIf
\end{algorithmic}
\end{algorithm}

Обратите внимание на асимметрию знака --- ею управляется всё содержимое
разд.~\ref{sec:num_pitfalls}: коагуляция есть $2 \to 1$ и, следовательно,
\emph{сток} числа частиц, фрагментация есть $1 \to 2$ и, следовательно,
\emph{источник}.

\section{Десять трюков}
\label{sec:num_tricks}

\subsection{Мажорантная частота столкновений}
\label{sec:trick1}

Вычислять точную парную скорость для каждой пары --- это $O(N^2)$ на шаг.
Вместо этого мажоранта $\fmaj \ge K$ управляет доминирующим пуассоновским
процессом, а события прореживаются с вероятностью принятия $K/\fmaj$.
Прореженный процесс точен: мажоранта влияет на стоимость, но никогда --- на
статистику.

\emph{Источники.} Ivanov \& Rogasinsky, Sov.\ J.\ Numer.\ Anal.\ Math.\
Modelling \textbf{3}, 453 (1988) [схема MCF в DSMC]; Garcia, van den
Broeck, Aertsens \& Serneels, Physica~A \textbf{143}, 535 (1987)
[мажоранта для коагуляции]; Gillespie, J.~Atmos.\ Sci.\ \textbf{32}, 1977
(1975) [точная стохастическая коалесценция].

\subsection{Бин-парная таблица мажоранты}
\label{sec:trick2}

Мажоранта --- это матрица $\fmaj_{bc}$ по парам массовых бинов, а не одно
глобальное число. Покупается этим \emph{стоимость} и \emph{проверяемость}.
Стоит проговорить, что именно, потому что корректность мажорантой купить
нельзя: она сокращается из физического темпа в точности.

\paragraph{Мажоранта сокращается.}
Тики доминирующего процесса приходят с темпом $wR/(2V)$; тик выбирает
упорядоченную пару $(i,j)$ с вероятностью $\fmaj_{ij}/R$; принимается он с
вероятностью $K_{ij}/\fmaj_{ij}$. Темп принятых событий на этой паре, стало
быть, равен
\begin{equation}
  \frac{wR}{2V}\cdot\frac{\fmaj_{ij}}{R}\cdot\frac{K_{ij}}{\fmaj_{ij}}
  \;=\; \frac{w\,K_{ij}}{2V},
  \label{eq:num_cancel}
\end{equation}
а суммирование по неупорядоченным парам даёт $w\sum_{i<j}K_{ij}/V$ ---
смолуховский темп, вообще не зависящий от $\fmaj$. Мажоранта сокращается
дважды: один раз в весе выбора пары, второй раз в доле принятия.
Следовательно, \emph{любая} допустимая мажоранта возвращает одну и ту же
физику; плохая лишь впустую тратит попытки. Равенство~\eqref{eq:num_cancel}
верно потому, что часы продвигаются на каждой \emph{попытке}, а не на каждом
событии --- см.\ разд.~\ref{sec:trick5}, где разобрана та единственная
правка, которая это ломает.

\paragraph{Что таблица покупает на самом деле.}
Одно глобальное $f_{\max} = K(m_{\max},m_{\max})$ --- вполне законная
мажоранта, дающая ровно ту же статистику. И тем не менее это неверный выбор,
по трём причинам:

\begin{itemize}
  \item \textbf{Доля принятия.} Типичная пара принимается с вероятностью
        $\langle K\rangle/f_{\max} \simeq (m_0/m_{\max})^{\lambda}$. На $D$
        декадах это $10^{-\lambda D}$: для геометрического ядра на четырёх
        декадах --- $2\times10^{-3}$, то есть около пятисот попыток на
        событие против $1{,}05$ с таблицей. Таблица же держит долю принятия
        $O(1)$ в \emph{каждой} паре бинов, потому что граница берётся в углах
        именно той пары, которую разыграли; перебор тогда составляет всего
        $10^{0.1\lambda}$ на аргумент при бинах в $0{,}1$~декады.
  \item \textbf{Ложное срабатывание сторожа зависания.}
        \code{max\_stall\_tries} считает попытки без принятого события. При
        доле принятия $2\times10^{-3}$ его умолчание $2\times10^{6}$ ---
        это всего $\sim\!4000$ событий честной работы, так что здоровый
        прогон может быть убит с диагнозом <<acceptance collapsed>>.
  \item \textbf{Отстающий максимум --- настоящее смещение.} \emph{Бегущий}
        глобальный максимум, не обновлённый до того, как какая-то частица
        переросла его, даёт $\fmaj < K$, и прореживание тогда молча неверно;
        это, в отличие от двух пунктов выше, уже отказ по корректности.
        Угловая таблица статична, и каждая попытка явно проверяет
        $K \le \fmaj$, поэтому та же ситуация приводит к падению, а не к
        смещению.
\end{itemize}

\emph{Источник.} Eibeck \& Wagner, SIAM J.~Sci.\ Comput.\ \textbf{22}, 802
(2000),
\url{https://doi.org/10.1137/S1064827599353488}.

\subsection{Членство в бине за $O(1)$}
\label{sec:trick3}

\code{bins[b]} --- список индексов частиц, а \code{pos\_in\_bin[i]} ---
позиция $i$ внутри него. Удаление переставляет последний элемент в
образовавшуюся дырку, поэтому вставка и удаление стоят $O(1)$ без
перевыделения памяти. Тот же приём swap-and-pop применяется и к самим
массивам частиц при удалении --- вот почему любой индекс, переживший
удаление, обязан быть перепроверен.

\subsection{Инкрементальное ведение скоростей}
\label{sec:trick4}

Когда заселённость одного бина меняется на $\delta$,
\[
  T \mathrel{+}= \delta\, \fmaj_{\cdot b},
  \qquad
  \text{row} \gets N \odot (T - \mathrm{diag}\,\fmaj),
  \qquad
  R \gets \textstyle\sum_b \text{row}_b ,
\]
что есть векторная операция $O(B)$ вместо пересчёта полного
матрично-векторного произведения. Выбор использует векторизованные
\code{cumsum} и \code{searchsorted}, а не цикл на Python; иначе именно этот
цикл --- доминирующая стоимость. При очень больших $B$ дерево Фенвика даёт
$O(\log B)$ (Goodson \& Kraft, J.~Comput.\ Phys.\ \textbf{183}, 210 (2002)).

\subsection{Часы}
\label{sec:trick5}

Шаг равен
\begin{equation}
  \Delta t \;=\; \frac{2V}{w\,R},
  \label{eq:num_clock}
\end{equation}
и три отдельные вещи в нём легко испортить.

\paragraph{Упорядоченные и неупорядоченные пары.}
$R$ считает каждую неупорядоченную пару дважды, уравнение~\eqref{eq:num_R},
тогда как физическая мажорантная скорость пробегает неупорядоченные пары.
Отсюда в числителе $2V$, а не $V$. На \emph{выбор} пары это не влияет ---
оба порядка ведут к одной и той же неупорядоченной паре, и множители
сокращаются, --- поэтому ошибка невидима в динамике, перемасштабирует $t$
ровно вдвое и ломает любое сравнение с аналитическим решением.

\paragraph{Веса.}
Здесь $w$ --- константа, но он всё равно вписан в шаг, чтобы формула в
точности совпадала со взвешенным вариантом и чтобы постоянное $w \neq 1$
просто перемасштабировало плотность. Ничто в этом варианте никогда не
меняет $w$, поэтому часы не могут быть испорчены учётом весов. Ради этого
базовый вариант и существует.

\paragraph{Квадратичность.}
$R \propto N^2$, потому что процесс парный, следовательно
\begin{equation}
  t \;\simeq\; \frac{2E}{N^2}
  \label{eq:num_t_of_E}
\end{equation}
после $E$ попыток. Шаг, линейный по $N$, описывает унарный процесс и даёт
неверное $N(t)$. У уравнения~\eqref{eq:num_t_of_E} есть практическое
следствие, на которое легко наткнуться случайно: \emph{увеличение $N$ при
фиксированном бюджете событий делает прогон короче по физическому времени,
а не длиннее}. Любые ворота, заданные абсолютным временем и настроенные при
одном $N$, при $10N$ не пересекаются никогда.

\paragraph{Время идёт на каждой попытке, а не на каждом событии.}
\code{t\_phys} увеличивается сразу после вычисления $\Delta t$, за полсотни
строк до проверки на принятие, так что нулевые (прореженные) попытки тоже
двигают часы. Именно это заставляет работать сокращение
\eqref{eq:num_cancel}, и перенос этого увеличения внутрь ветки принятия ---
та единственная правка, которая превращает мажоранту из стоимости в
\emph{смещение}. Физическое время продвигалось бы тогда на $2V/(wR)$ за одно
принятое событие, поэтому темп событий вышел бы равным $R$ вместо
$R\,\langle K/\fmaj\rangle = \sum K$, и при глобальном $f_{\max}$
\[
  \Delta t_{\rm эфф} \;\propto\; \frac{1}{N^2 f_{\max}}
                     \;\propto\; \frac{1}{N^2 m_{\max}^{\lambda}} .
\]
Ядро вошло бы в часы через \emph{экстремум} распределения, а не через
типичную массу. В открытой системе $m_{\max}$ прижат к стоку и не растёт,
так что зависимость от $m_0$ ушла бы из часов целиком, дрейф выродился бы в
$dm_0/dt \propto m_0$ (экспоненциальный, $\beta = 1$), и условие постоянного
потока вернуло бы $\alpha = -2$ при любом ядре. Это третий, чисто численный
путь к тому же отпечатку $-2$, который разд.~\ref{sec:trick8} производит
физически.

\subsection{Детерминированная инжекция по остатку}
\label{sec:trick6}

Инжекция использует дробный накопитель,
\[
  \nu \mathrel{+}= q\,\Delta t/w,
  \qquad k = \lfloor \nu \rfloor,
  \qquad \nu \mathrel{-}= k,
\]
а не пуассоновский розыгрыш. Это убирает дробовой шум на масштабе инжекции.
Это модельный выбор, и он записывается в метаданные вывода, а не остаётся
неявным.

\subsection{Учёт изохрон}
\label{sec:trick7}

Каждая частица несёт время своего рождения; её возраст есть
$\tau = t - t_{\rm born}$. Снапшоты накапливают двумерную гистограмму по
$(\tau, m)$, строки которой суть изохроны, а сумма по строкам ---
усреднённый по времени спектр.

Для \textbf{фрагментации} возраст осколка однозначен: осколки происходят
ровно от одного родителя. Для \textbf{коагуляции} у слившейся частицы два
предка, и правило наследования приходится выбирать. Оно вынесено в
параметр \code{age\_rule} именно потому, что анализ изохрон от него
зависит, а сделанный выбор обязан быть заявлен и проверен:
\[
  t_{\rm born} \;=\;
  \begin{cases}
    \min(t_1,t_2)                    & \text{\code{'min'}: побеждает старший предок,}\\[2pt]
    t_1 \ \text{если}\ m_1 \ge m_2   & \text{\code{'heavier'},}\\[2pt]
    \dfrac{m_1 t_1 + m_2 t_2}{m_1+m_2} & \text{\code{'mass\_weighted'}.}
  \end{cases}
\]

\subsection{Локальность правила фрагментации --- ловушка физики}
\label{sec:trick8}

Если правило звучит как <<ломается более тяжёлая из пары, чем бы в неё ни
попали>>, то скорость разрушения тела массы $m_0$ на степенном фоне
$F(m) \propto m^{\alpha}$ равна
\begin{equation}
  \nu(m_0) \;\sim\; \int_{\msink}^{m_0} F(m')\,K(m_0,m')\,dm'
           \;\sim\; \int_{\msink/m_0} x^{\alpha}\,dx ,
  \qquad x = m'/m_0 ,
\end{equation}
и этот интеграл \emph{расходится на нижнем пределе всякий раз, когда
$\alpha < -1$}. Дрейф тогда задаётся самыми мелкими частицами в ящике ---
масштабом стока, а не $m_0$. Среднеполевое замыкание теряет силу, дрейф
вырождается в $dm_0/d\tau \sim -C m_0$ (то есть $\beta = 1$,
экспоненциальный спад), и измеряемый индекс запирается на $-2$ независимо
от ядра: мнимая универсальность, которая на деле есть граничный эффект.

Лекарство --- то же, что заложил Dohnanyi (1969): порог катастрофического
разрушения. Пылинка не разбивает валун. Требование
\begin{equation}
  m_{\rm small} \;\ge\; f\,m_{\rm large},
  \qquad f = O(0{,}1\text{--}1),
\end{equation}
ограничивает интеграл областью $x \in [f,1]$, которая свободна от масштаба,
восстанавливает локальность и возвращает $\alpha = -(3+\lambda)/2$.

\emph{Замкнутая} система к этому невосприимчива: её автомодельный пакет
узок, так что все партнёры уже порядка $m_0$ и правило локально по
построению. Именно поэтому замкнутые прогоны согласуются с теорией при
$f=0$, а открытые --- нет.

\subsection{Отсутствие ресэмплинга --- и его цена}
\label{sec:num_cost}

Что покупается:
\begin{itemize}
  \item $\sum w$ и $\sum wm$ сохраняются с машинной точностью; выводимый
        \code{mass\_drift} тождественно нулевой, и любое ненулевое значение ---
        настоящая ошибка;
  \item никаких тонкостей с часами вообще, $\Delta t = 2V/R$ при $w=1$;
  \item ничто не может сместить результат, потому что ничто не
        пересэмплируется.
\end{itemize}

Что платится, и это связывающее ограничение для каждого прогона:
\begin{itemize}
  \item \textbf{Коагуляция.} В замкнутом ящике $N = N_i m_i/m_0$, поэтому
        чтобы наверху оставалось $\sim100$ живых частиц, требуется
        $m_0 \lesssim N_i/100$: при $N_i = 10^6$ это четыре декады, при
        $N_i = 10^8$ --- шесть (и $\sim3$~ГБ массивов). Если продавить
        дальше, прогон заканчивается горсткой частиц; при
        $N_i = 2\times10^5$ один тест финишировал с единственной живой
        частицей, а полезное плато сжалось с $5{,}6$ декад до $2{,}7$.
  \item \textbf{Фрагментация.} Хуже, и в другую сторону: $N$ растёт как
        $m_i/m_0$, так что завершённый каскад держит $N_0\,\minj/\msink$
        частиц. Три декады умножают популяцию в тысячу раз, и \emph{прогон
        заканчивает память, а не физика}.
\end{itemize}

Пользуйтесь этим вариантом, когда приоритет --- безупречная опорная точка.
Пользуйтесь взвешенным вариантом, добавляющим размножение и прореживание
популяции, когда приоритет --- диапазон: десять декад там стоили
$4{,}5\times10^5$ событий и дали $\alpha = -1{,}005$ против теоретического
$-1{,}000$. Оба согласуются всюду, где могут работать оба, --- ради этого
они и держатся вместе.

\subsection{Размещение снапшотов}
\label{sec:trick10}

Наблюдаемая в замкнутой системе --- это сумма Римана
$F(m) = \sum_k (dN/dm)(m,t_k)\,\Delta t_k$, и у выборки по $t$ есть две
раздельные работы, которые нельзя смешивать. \textbf{Размещение} решает,
что будет \emph{разрешено}: при данном $m$ по\-дын\-те\-граль\-ное выражение
имеет
пик при $m_0 \sim m$, поэтому каждой декаде $m$ нужны снапшоты в
соответствующей декаде $m_0$. \textbf{Вес} решает, что будет
\emph{вычислено}: он обязан быть фактическим $\Delta t$ между соседними
снапшотами, каким бы ни было размещение.

\begin{center}
\footnotesize
\begin{tabular}{@{}l p{6.6cm} l@{}}
\toprule
размещение & что происходит & измеренное $\alpha$ \\
 & & (теория $-1$) \\
\midrule
\code{events} & число событий на единицу $m_0$ падает как $m_0^{-2}$,
  поэтому $\Delta t$ взрывается; $94\%$ всего веса приходится на один
  снапшот & $-1{,}4 \ldots -2{,}6$ \\
\code{t\_phys}, узкий диапазон & хорошо работает на $\lesssim 2$ декадах
  & $-0{,}97 \ldots -1{,}02$ \\
\code{t\_phys}, широкий & $t \propto m_0^{1-\lambda}$, поэтому равномерность
  по $t$ есть почти равномерность по $m_0$; на десяти декадах нижние пять
  не получают ни одной выборки & $-0{,}04 \pm 0{,}16$ \\
\code{log\_m0} & $k$ выборок на декаду $m_0$, везде
  & $\mathbf{-1{,}005}$, $R^2 = 0{,}9995$ \\
\bottomrule
\end{tabular}
\end{center}

Провал режима \code{events} заслуживает отдельного упоминания, потому что
он бесшумен: при таком размещении искатель плато из
разд.~\ref{sec:num_measure} всё равно сообщает о чистом плато шириной в
декаду со значением $-1{,}375 \pm 0{,}14$. Внутренне согласованное,
полностью неверное измерение. \emph{Подгонка по плато не спасает сломанный
оценщик}, потому что она проверяет, есть ли у вычисленной вами кривой
скейлинговая область, а не является ли эта кривая правильной величиной.
Независимая гарантия --- печатаемый вес последнего снапшота: он обязан быть
$\ll 1\%$.

Ресэмплинг (разд.~\ref{sec:num_cost}) и размещение лечат разные болезни.
Ресэмплинг покупает \emph{диапазон}; размещение покупает \emph{корректное
измерение} на этом диапазоне.

\paragraph{Правило.} Замкнутая система $\Rightarrow$ \code{log\_m0}, шаг $=$
число снапшотов на декаду ($20$--$30$). Открытая система $\Rightarrow$
стационарное состояние есть \emph{последний} снапшот, $\Delta t$ в оценщик
не входит вовсе, и годится любое расписание.

\section{Сток как критерий}
\label{sec:num_sink}

Критерий остановки или открытия ворот, сформулированный в событиях, есть
утверждение о \emph{ресурсах}; сформулированный в поглощениях на стоке ---
утверждение о \emph{физике}. Только второй переживает смену $N$, диапазона
масс или ядра. Поэтому движок считает $n_{\rm out}$ --- число физических
частиц, покинувших систему через поглощающую границу, --- и выставляет его
наружу тремя способами.

\begin{center}
\small
\begin{tabular}{@{}l p{8.2cm}@{}}
\toprule
\code{iso\_start\_sink} (по умолчанию $10$) &
  начинать накопление изохрон только после $n_{\rm out}$ поглощений \\
\code{stop\_sink\_events} (по умолчанию \code{None}) &
  останавливать прогон на $n_{\rm out}$ поглощениях \\
\code{out["n\_out"]} &
  сам счётчик как временной ряд по снапшотам \\
\bottomrule
\end{tabular}
\end{center}

Ворота говорят: не биновать ни одного возраста, пока каскад доказуемо не
протолкнул столько-то частиц через \emph{весь} инерционный интервал. Их
стоимость следует из баланса массы --- ящик держит $\simeq N\sqrt{\msink}$ в
стационаре, а $n_{\rm out}$ поглощений уносят $n_{\rm out}\msink$, --- так
что
\begin{equation}
  E \;\simeq\; N\sqrt{\msink} \;+\; n_{\rm out}\,\msink .
  \label{eq:num_budget}
\end{equation}
Второе слагаемое не зависит от $N$: ворота не дорожают с ростом прогона,
дорожает только заполнение ящика. Уравнение~\eqref{eq:num_budget} выведено
при $\lambda = 0$ и переоценивает стоимость при $\lambda > 0$, где более
крутое ядро гонит массу к стоку быстрее в расчёте на событие.

У замкнутой системы стока нет, поэтому взведённое значение по умолчанию там
молча разоружается; исключение бросается только на явно переданное
значение. Различие между <<это набрал пользователь>> и <<это пришло из
сигнатуры>> несёт подкласс \code{int}, так что во всех прочих местах
значение ведёт себя как обычное целое.

\section{Измерительный слой}
\label{sec:num_measure}

Прогон производит массивы; превращение их в показатель степени --- отдельная
стадия со своими режимами отказа.

\paragraph{Оценщик.} Замкнутая система: про\-ин\-те\-гри\-ро\-ван\-ная по
возрасту суперпозиция $F(m) = \sum_t (dN/dm)\,\Delta t$. Открытая система:
стационарное состояние \emph{и есть} мгновенный спектр, поэтому оценщиком
служит последний снапшот, а $\Delta t$ в него не входит.

\paragraph{Локальный наклон и плато.} Честный способ смотреть на измеренный
спектр --- это
\begin{equation}
  \Gamma(m) \;=\; \frac{d\log F}{d\log m},
\end{equation}
вычисленный методом наименьших квадратов на скользящем окне. Настоящий
инерционный интервал --- это \emph{плато} в $\Gamma$; загрязнённые концы
проявляются как места, где $\Gamma$ загибается. Сообщаемый индекс есть
среднее $\Gamma$ по самому длинному непрерывному участку, на котором оно
остаётся плоским в пределах допуска, вместе с его разбросом и шириной в
декадах. Плато уже примерно одной декады не является степенным законом,
сколь бы узкой ни была его погрешность; ширина печатается рядом с индексом
именно поэтому. Если ни один участок не проходит, ответ --- \code{nan}, и
это правильный ответ, а не отказ.

\paragraph{Защитная полоса (guard band).} Независимо от этого: отрезать по
$0{,}5$~декады с каждого конца интервала между двумя характерными массами
задачи --- $[\minj,\msink]$ для открытой системы, $[m_i, m_0^{\max}]$ для
замкнутой --- и подогнать там единственный степенной закон. Эта полоса
выбирается \emph{до} того, как посмотрели на данные. Две оценки обязаны
согласоваться; если нет, каскад не развился на том диапазоне, который был
предположен.

Подгонка по всему массиву не даёт ни того, ни другого. На одном прогоне
замкнутой коагуляции при $\lambda=0$ и теоретическом $\alpha=-1$: весь
массив $-1{,}71$, защитная полоса $-1{,}09$, плато $-1{,}16 \pm 0{,}15$ на
$1{,}3$ декады.

\paragraph{Компенсированное представление.} На панелях рисуется
\begin{equation}
  m^2 \frac{dN}{dm} \;=\; \text{масса на логарифмический интервал массы},
\end{equation}
чей наклон равен $\alpha+2$. Вся подгонка делается по сырому $dN/dm$;
компенсация применяется только при отрисовке. Модельная линия несёт индекс
плато и амплитуду, полученную как свободный член наименьших квадратов при
\emph{фиксированном} наклоне, так что никакая вторая подгонка не может
тихо сдвинуть показатель, а теоретическая линия привязывается тем же
способом --- тогда две кривые различаются только наклоном.

\paragraph{Предсказание.} При степени однородности ядра $\lambda$
\begin{equation}
  \beta =
    \begin{cases}
      \lambda & \text{замкнутая: пакет } n \propto m_0^{-2},\\[2pt]
      \dfrac{1+\lambda}{2} & \text{открытая: фон } n \propto m_0^{\alpha},
    \end{cases}
  \qquad
  \alpha = -(1+\beta),
  \quad
  b = \frac{1}{1-\beta}.
\end{equation}
Коагуляция и фрагментация разделяют $\beta$, $b$ и $\alpha$ в точности;
переворачивается только знак дрейфа, поэтому закон читается в $\tau$ при
движении вверх и в $\tau_\ast - \tau$ при движении вниз.

\begin{center}
\footnotesize
\begin{tabular}{@{}@{}l@{\ \ }l@{\ \ }c@{\ \ }c@{\ \ }c@{}}
\toprule
ядро & $K(m_1,m_2)$ & $\lambda$ &
  замк.\ $\alpha$, $b$ &
  откр.\ $\alpha$, $b$ \\
\midrule
постоянное     & $1$                                & $0$   & $-1$,\ \ $1$    & $-3/2$,\ \ $2$ \\
сумма          & $m_1^{1/3}+m_2^{1/3}$              & $1/3$ & $-4/3$,\ \ $3/2$& $-5/3$,\ \ $3$ \\
геометрич.\    & $(m_1^{1/3}+m_2^{1/3})^2$          & $2/3$ & $-5/3$,\ \ $3$  & $-11/6$,\ \ $6$ \\
аддитивное     & $m_1+m_2$                          & $1$   & $-2$,\ \ $\infty$ & $-2$,\ \ $\infty$ \\
произвед.\     & $m_1 m_2$                          & $2$   & \multicolumn{2}{c}{выше гелеобразования} \\
\bottomrule
\end{tabular}
\end{center}

\section{Валидация}
\label{sec:num_validation}

Постоянное ядро, $\lambda = 0$. \emph{плато} --- автоматический индекс,
\emph{полоса} --- априорная защитная полоса, \emph{дек} --- ширина плато в
декадах.

\begin{center}
\small
\begin{tabular}{@{}l r r r r r r@{}}
\toprule
случай & $b$ & $b_{\rm теор}$ & плато & $\pm$ & дек & полоса \\
\midrule
1 замкнутая коагуляция     & $0{,}990$ & $1$ & $-1{,}077$ & $0{,}126$ & $2{,}70$ & $-1{,}035$ \\
2 замкнутая фрагментация   & $1{,}001$ & $1$ & $-0{,}978$ & $0{,}141$ & $3{,}00$ & $-1{,}056$ \\
3 открытая коагуляция      & $2{,}001$ & $2$ & $-1{,}475$ & $0{,}095$ & $3{,}60$ & $-1{,}498$ \\
4 открытая фрагм., $f=0$   & ---       & $2$ & $-2{,}047$ & $0{,}183$ & $0{,}80$ & $-2{,}354$ \\
4 открытая фрагм., $f=0{,}3$ & ---     & $2$ & $-1{,}550$ & $0{,}091$ & $3{,}40$ & $-1{,}517$ \\
\bottomrule
\end{tabular}
\end{center}

Случай~1 дополнительно воспроизводит точное решение Смолуховского для
постоянного ядра, $m_0(t) = m_i(1+n_0K_1t/2)$, с точностью
$\max|{\rm отношение}-1| = 7{,}8\times10^{-14}$ --- прямая проверка
уравнения~\eqref{eq:num_clock}. Дрейф массы тождественно нулевой во всех
четырёх случаях.

Две строки следует читать как диагноз, а не как измерение. Случай~4 при
$f=0$ возвращает $-2{,}05$ на $0{,}80$ декады: индекс залип на $-2$ в
точности так, как предсказывает разд.~\ref{sec:trick8}, а узкое плато само
объявляет, что это не измерение. И показатель роста $b$ для случая~4
отсутствует потому, что при таком бюджете более $99\%$ живых частиц всё ещё
происходят от начального условия и наследуют $t_{\rm born}=0$; ось возраста
тогда вырождена --- возраст равен астрономическому времени сразу у всех, ---
и изохроны деградируют в <<спектр как функция времени>>. Тот же факт
проявляется как $M_{\rm out}/M_{\rm in} > 1$: сток сливает начальный
резервуар, чья масса в $M_{\rm in}$ никогда не входила.

\section{Практические режимы отказа}
\label{sec:num_pitfalls}

\paragraph{Масса вне сетки --- жёсткое падение.}
Мажоранта строится по \emph{верхним углам} бинов, а $b(m)$ зажимает всё,
что выше верхней границы, в последний бин, чей угол тогда меньше самой
частицы. Следовательно $K(m,m) > \fmaj_{B-1,B-1}$, и первая же пара бросает
\code{Majorant violated}. При постоянном ядре этого не может случиться в
принципе, поскольку $K\equiv1$ плоское; при любом $\lambda>0$ это
происходит немедленно. Проверяйте, что сетка охватывает $\minj$ и $\msink$,
до запуска, а не после.

\paragraph{Неограниченный бюджет событий безопасен ровно в одном случае.}
Снимать ограничение на число событий можно только тогда, когда прогон
способно завершить что-то другое:

\begin{center}
\small
\begin{tabular}{@{}l l p{7.4cm}@{}}
\toprule
случай & без предела & что останавливает --- или не останавливает \\
\midrule
замкнутая коагуляция   & безопасно & \code{stop\_max\_mass}; к тому же $N$ монотонно
                                     убывает, так что \code{live}~$<2$ --- гарантированный
                                     терминатор \\
замкнутая фрагментация & опасно    & $N$ растёт; сторож зависания срабатывает лишь
                                     тогда, когда \emph{всё} окажется ниже пассивного
                                     пола, то есть при $N_0 m_i/m_{\rm floor}$ частицах \\
открытая коагуляция    & зависнет  & $N$ сидит на полке, поэтому никакое другое
                                     условие сработать не может \\
открытая фрагментация  & опасно    & $N$ растёт неограниченно \\
\bottomrule
\end{tabular}
\end{center}

В открытых случаях правильный тормоз --- \code{stop\_sink\_events}
(разд.~\ref{sec:num_sink}); дополнительно оставленный ресурсный потолок не
стоит ничего и ограничивает последствия неверно настроенной инжекции.
Записанная причина остановки затем говорит, какой из двух сработал.

\paragraph{Смена ядра обесценивает четыре константы.}
Однородность $\lambda$ и все предсказанные показатели подхватываются
автоматически, а искатель плато, защитная полоса и компенсация вообще не
знают про ядро. Четыре вещи --- знают:

\begin{enumerate}
  \item \textbf{Темп инжекции.} Баланс числа в стационаре есть
        $q = \langle K\rangle N^2/(2V)$. Запись $q = N^2/2$ --- это та же
        формула при $\langle K\rangle = 1$. Поскольку $N$ набирается на
        масштабе инжекции при $\alpha<-1$, большинство пар суть
        $(\minj,\minj)$, так что $\langle K \rangle \sim K(\minj,\minj)$, а
        более тяжёлые партнёры поднимают эту величину примерно вдвое.
  \item \textbf{Время столкновения.} $t_c = 1/(Kn)$, а не $1/n$.
  \item \textbf{Биновка по возрасту.} Один возрастной бин шириной
        $\Delta_\tau$ декад покрывает $b\,\Delta_\tau$ декад \emph{массы}.
        Фиксируйте разрешение по массе и выводите шаг по возрасту из него,
        $\Delta_\tau = 0{,}30/b$; иначе при $b=6$ шаг, настроенный под
        $b=2$, перепрыгивает почти декаду массы за бин, и двухдекадное окно
        подгонки ловит два бина.
  \item \textbf{Точные решения.} Замкнутая форма Смолуховского и проверка
        <<$m_0(t)$ обязано быть линейным>> --- свойства исключительно
        постоянного ядра.
\end{enumerate}

\paragraph{У отпечатка $-2$ три разные причины.}
Индекс $-2$ возвращают: (i) настоящее ядро с $\lambda = 1$, где это
правильный ответ; (ii) нелокальное правило разрушения из
разд.~\ref{sec:trick8}, где это граничный эффект; (iii) часы, продвигаемые
на событие, а не на попытку, разд.~\ref{sec:trick5}, где это численный
артефакт. Все три гонят $\beta \to 1$ и на бумаге выглядят одинаково. Само
по себе значение $-2$ поэтому не говорит ни о чём. Различающая проверка ---
меняется ли индекс при смене ядра: если не меняется, это не универсальность.

\paragraph{Плоская полка --- не инерционный интервал.}
В замкнутой фрагментации пассивный пол замораживает всё, что ниже него: эти
частицы больше никогда не дробятся и накапливаются в плоскую полку, к
каскаду отношения не имеющую. Оставленная в массиве, эта полка --- самый
длинный плоский участок из всех, $4{,}3$ декады при $\alpha \simeq 0$, --- и
искатель плато послушно сообщает именно о ней. Ограничивайте поиск областью
$m \ge m_{\rm floor}$.

% ---------------------------------------- удалить при вставке в main.tex
\end{document}

%      рядом с уже имеющимся \include{Parts/unified-cascade_v4}.
%      В преамбулу main.tex тогда нужно добавить:
%          \usepackage{algorithm}
%          \usepackage{algpseudocode}
%          \usepackage{booktabs}
%      и, поскольку main.tex сейчас английский, переключить babel на
%          \usepackage[english,russian]{babel}
%      и кодировку на \usepackage[T1,T2A]{fontenc}.
% =====================================================================

% ------------------------------------------ ПРЕАМБУЛА -- ОТРЕЗАТЬ ОТСЮДА
\documentclass[a4paper,12pt]{article}

\usepackage[T1,T2A]{fontenc}
\usepackage[utf8]{inputenc}
\usepackage[english,russian]{babel}
\usepackage{amsmath,amssymb}
\usepackage{geometry}
\geometry{margin=2.5cm}
\usepackage{graphicx}
\usepackage{booktabs}
\usepackage{algorithm}
\usepackage{algpseudocode}
\usepackage[unicode,colorlinks=true,linkcolor=blue,citecolor=blue,urlcolor=blue]{hyperref}

\newcommand{\minj}{m_{\mathrm{inj}}}
\newcommand{\msink}{m_{\mathrm{sink}}}
\newcommand{\fmaj}{f^{\mathrm{maj}}}
\newcommand{\code}[1]{\texttt{#1}}
\emergencystretch=2em   % русские слова длиннее английских

\begin{document}
\selectlanguage{russian}

\title{Численная схема: метод Монте-Карло с мажорантной частотой\\ для каскада в пространстве масс}
\author{}
\date{}
\maketitle
% ------------------------------------------ ПРЕАМБУЛА -- ОТРЕЗАТЬ ДОСЮДА

\section{Область применимости и принцип построения}
\label{sec:num_scope}

Движок интегрирует стохастический каскад в пространстве масс в четырёх
конфигурациях,
\[
  \text{процесс} \in \{\text{коагуляция},\ \text{фрагментация}\},
  \qquad
  \text{система} \in \{\text{замкнутая},\ \text{открытая}\},
\]
одним и тем же путём в коде. Между четырьмя случаями различаются ровно две
вещи:

\begin{enumerate}
  \item что принятое событие делает с двумя выбранными частицами, и
  \item граничные условия --- инжекция плюс поглощающий сток, либо ни того,
        ни другого.
\end{enumerate}

Всё остальное --- построение мажоранты, выбор пары, часы, учёт ---
общее. Это сознательное ограничение, а не удобство: расхождение между
двумя из четырёх случаев тогда нельзя списать на два разных интегратора.

Описываемый здесь вариант не несёт весов частиц. Одна расчётная частица
есть одна физическая частица, $w \equiv 1$ на весь прогон, поэтому
$\sum_i w_i$ и $\sum_i w_i m_i$ сохраняются с машинной точностью на каждом
шаге. Измеряемый дрейф массы тождественно равен нулю, и любое ненулевое
значение --- это ошибка в коде, а не статистическое отклонение.
Раздел~\ref{sec:num_cost} говорит, чего это стоит.

\subsection{Состояние}

Состояние --- четыре параллельных массива по живым частицам плюс разбиение
пространства масс на бины:

\begin{center}
\small
\begin{tabular}{@{}ll@{}}
\toprule
\code{mass[i]}       & масса частицы $m_i$ \\
\code{inj\_time[i]}  & время рождения, нужно для изохрон (разд.~\ref{sec:trick7}) \\
\code{bin\_of[i]}    & номер массового бина, в котором лежит $i$ \\
\code{pos\_in\_bin[i]} & позиция $i$ внутри списка индексов этого бина \\
\midrule
\code{edges[0..B]}   & $B+1$ логарифмически расположенных границ бинов \\
\code{bins[b]}       & список индексов частиц в бине $b$ \\
\code{N\_bin[b]}     & заселённость $N_b$ \\
\bottomrule
\end{tabular}
\end{center}

Частица массы $m$ попадает в бин
$b(m) = \max\bigl(0,\min(B-1,\ \mathrm{searchsorted}(\text{edges},m)-1)\bigr)$.
Зажим (clamping) в этом выражении --- ловушка, и она разобрана в
разд.~\ref{sec:num_pitfalls}.

\section{Базовый алгоритм}
\label{sec:num_algorithm}

\subsection{Инициализация}

\begin{algorithm}[h]
\caption{Инициализация: бины, таблица мажоранты, вектор скоростей}
\label{alg:setup}
\begin{algorithmic}[1]
\State разложить все начальные частицы по бинам; сформировать заселённости $N_b$
\State \textbf{таблица мажоранты:} для всех $b,c$ \quad
       $\fmaj_{bc} \gets K(\text{edges}[b+1],\ \text{edges}[c+1])$
       \Comment{верхние углы}
\State $T_b \gets \sum_c \fmaj_{bc} N_c$
\State $\text{row}_b \gets N_b\,(T_b - \fmaj_{bb})$
\State $R \gets \sum_b \text{row}_b$
\end{algorithmic}
\end{algorithm}

Угловое правило $\fmaj_{bc}=K(\text{edges}[b+1],\text{edges}[c+1])$
ограничивает $K$ сверху на всей паре бинов \emph{при условии, что $K$ не
убывает по обоим аргументам}. Это единственное структурное требование к
ядру, и оно проверяется во время счёта, а не предполагается: каждая
принятая попытка проверяет $K(m_i,m_j) \le \fmaj_{bc}$ и падает, если это
нарушено.

По построению
\begin{equation}
  R \;=\; \sum_b N_b\bigl(T_b - \fmaj_{bb}\bigr)
    \;=\; \sum_{i \neq j} \fmaj_{ij}
    \;=\; 2 \sum_{i<j} \fmaj_{ij},
  \label{eq:num_R}
\end{equation}
то есть $R$ пробегает \emph{упорядоченные} пары. Двойка в
уравнении~\eqref{eq:num_R} --- источник двойки в часах,
уравнение~\eqref{eq:num_clock}.

\subsection{Главный цикл}

\begin{algorithm}[h]
\caption{Главный цикл: одна мажорантная попытка}
\label{alg:mainloop}
\begin{algorithmic}[1]
\While{ни одно условие остановки не выполнено}
  \State $\Delta t \gets 2V/(wR)$; \quad $t \gets t + \Delta t$
         \Comment{часы, разд.~\ref{sec:trick5}}
  \State \textbf{инжекция (открытые системы):}
         $\nu \gets \nu + q\,\Delta t / w$; \quad
         $k \gets \lfloor \nu \rfloor$; \quad $\nu \gets \nu - k$;
         добавить $k$ частиц массы $\minj$
  \State выбрать бин-строку $b$ с $P(b) \propto \text{row}_b$
  \State выбрать бин-столбец $c$ с
         $P(c) \propto N_c \fmaj_{bc}$, беря $(N_b-1)\fmaj_{bb}$ при $c=b$
  \State выбрать $i$ равномерно из \code{bins[b]}, $j \neq i$ равномерно из \code{bins[c]}
  \State \textbf{прореживание:} принять с вероятностью $K(m_i,m_j)/\fmaj_{bc}$;
         иначе \textbf{continue} \Comment{нулевое событие}
  \State применить событие (алг.~\ref{alg:event}); применить поглощающий сток
  \State записать снапшот, если этого требует расписание (разд.~\ref{sec:trick10})
\EndWhile
\end{algorithmic}
\end{algorithm}

Отклонённая попытка --- не ошибка: именно прореживание делает доминирующий
процесс точным. Доля принятия, выводимая как \code{acc} в логе прогона,
есть $O(1)$ по построению бин-парной таблицы --- измерено $0{,}95$ для
геометрического ядра на бинах шириной $0{,}1$~декады, где угловой перебор
составляет $10^{0.1\lambda}$ на аргумент.

\subsection{Событие}

\begin{algorithm}[h]
\caption{Принятое событие --- единственное место, где четыре случая расходятся}
\label{alg:event}
\begin{algorithmic}[1]
\If{коагуляция}
  \State $m_{\mathrm{new}} \gets m_i + m_j$
  \State унаследовать возраст $i$ от пары $(i,j)$ по \code{age\_rule} (разд.~\ref{sec:trick7})
  \State $m_i \gets m_{\mathrm{new}}$; перебиновать $i$; \textbf{удалить} $j$
  \State \textbf{если} открытая \textbf{и} $m_{\mathrm{new}} \ge \msink$ \textbf{то}
         удалить $i$; $M_{\mathrm{out}} \mathrel{+}= w\,m_{\mathrm{new}}$;
         $n_{\mathrm{out}} \mathrel{+}= 1$
\Else \Comment{фрагментация}
  \State $p \gets \arg\max(m_i,m_j)$; \quad $s \gets \arg\min(m_i,m_j)$
  \State \textbf{если} $m_s < f\,m_p$ \textbf{то continue}
         \Comment{порог разрушения, разд.~\ref{sec:trick8}}
  \State $\xi \sim U(0,1)$; \quad $m_a \gets \xi m_p$, \ $m_b \gets (1-\xi)m_p$
  \State оба осколка наследуют $t_{\mathrm{born}}$ единственного родителя $p$
  \State $m_p \gets m_a$; перебиновать $p$; \textbf{добавить} частицу массы $m_b$
  \State \textbf{если} открытая: удалить каждый осколок с $m \le \msink$,
         начиная со старшего индекса
\EndIf
\end{algorithmic}
\end{algorithm}

Обратите внимание на асимметрию знака --- ею управляется всё содержимое
разд.~\ref{sec:num_pitfalls}: коагуляция есть $2 \to 1$ и, следовательно,
\emph{сток} числа частиц, фрагментация есть $1 \to 2$ и, следовательно,
\emph{источник}.

\section{Десять трюков}
\label{sec:num_tricks}

\subsection{Мажорантная частота столкновений}
\label{sec:trick1}

Вычислять точную парную скорость для каждой пары --- это $O(N^2)$ на шаг.
Вместо этого мажоранта $\fmaj \ge K$ управляет доминирующим пуассоновским
процессом, а события прореживаются с вероятностью принятия $K/\fmaj$.
Прореженный процесс точен: мажоранта влияет на стоимость, но никогда --- на
статистику.

\emph{Источники.} Ivanov \& Rogasinsky, Sov.\ J.\ Numer.\ Anal.\ Math.\
Modelling \textbf{3}, 453 (1988) [схема MCF в DSMC]; Garcia, van den
Broeck, Aertsens \& Serneels, Physica~A \textbf{143}, 535 (1987)
[мажоранта для коагуляции]; Gillespie, J.~Atmos.\ Sci.\ \textbf{32}, 1977
(1975) [точная стохастическая коалесценция].

\subsection{Бин-парная таблица мажоранты}
\label{sec:trick2}

Мажоранта --- это матрица $\fmaj_{bc}$ по парам массовых бинов, а не одно
глобальное число. Покупается этим \emph{стоимость} и \emph{проверяемость}.
Стоит проговорить, что именно, потому что корректность мажорантой купить
нельзя: она сокращается из физического темпа в точности.

\paragraph{Мажоранта сокращается.}
Тики доминирующего процесса приходят с темпом $wR/(2V)$; тик выбирает
упорядоченную пару $(i,j)$ с вероятностью $\fmaj_{ij}/R$; принимается он с
вероятностью $K_{ij}/\fmaj_{ij}$. Темп принятых событий на этой паре, стало
быть, равен
\begin{equation}
  \frac{wR}{2V}\cdot\frac{\fmaj_{ij}}{R}\cdot\frac{K_{ij}}{\fmaj_{ij}}
  \;=\; \frac{w\,K_{ij}}{2V},
  \label{eq:num_cancel}
\end{equation}
а суммирование по неупорядоченным парам даёт $w\sum_{i<j}K_{ij}/V$ ---
смолуховский темп, вообще не зависящий от $\fmaj$. Мажоранта сокращается
дважды: один раз в весе выбора пары, второй раз в доле принятия.
Следовательно, \emph{любая} допустимая мажоранта возвращает одну и ту же
физику; плохая лишь впустую тратит попытки. Равенство~\eqref{eq:num_cancel}
верно потому, что часы продвигаются на каждой \emph{попытке}, а не на каждом
событии --- см.\ разд.~\ref{sec:trick5}, где разобрана та единственная
правка, которая это ломает.

\paragraph{Что таблица покупает на самом деле.}
Одно глобальное $f_{\max} = K(m_{\max},m_{\max})$ --- вполне законная
мажоранта, дающая ровно ту же статистику. И тем не менее это неверный выбор,
по трём причинам:

\begin{itemize}
  \item \textbf{Доля принятия.} Типичная пара принимается с вероятностью
        $\langle K\rangle/f_{\max} \simeq (m_0/m_{\max})^{\lambda}$. На $D$
        декадах это $10^{-\lambda D}$: для геометрического ядра на четырёх
        декадах --- $2\times10^{-3}$, то есть около пятисот попыток на
        событие против $1{,}05$ с таблицей. Таблица же держит долю принятия
        $O(1)$ в \emph{каждой} паре бинов, потому что граница берётся в углах
        именно той пары, которую разыграли; перебор тогда составляет всего
        $10^{0.1\lambda}$ на аргумент при бинах в $0{,}1$~декады.
  \item \textbf{Ложное срабатывание сторожа зависания.}
        \code{max\_stall\_tries} считает попытки без принятого события. При
        доле принятия $2\times10^{-3}$ его умолчание $2\times10^{6}$ ---
        это всего $\sim\!4000$ событий честной работы, так что здоровый
        прогон может быть убит с диагнозом <<acceptance collapsed>>.
  \item \textbf{Отстающий максимум --- настоящее смещение.} \emph{Бегущий}
        глобальный максимум, не обновлённый до того, как какая-то частица
        переросла его, даёт $\fmaj < K$, и прореживание тогда молча неверно;
        это, в отличие от двух пунктов выше, уже отказ по корректности.
        Угловая таблица статична, и каждая попытка явно проверяет
        $K \le \fmaj$, поэтому та же ситуация приводит к падению, а не к
        смещению.
\end{itemize}

\emph{Источник.} Eibeck \& Wagner, SIAM J.~Sci.\ Comput.\ \textbf{22}, 802
(2000),
\url{https://doi.org/10.1137/S1064827599353488}.

\subsection{Членство в бине за $O(1)$}
\label{sec:trick3}

\code{bins[b]} --- список индексов частиц, а \code{pos\_in\_bin[i]} ---
позиция $i$ внутри него. Удаление переставляет последний элемент в
образовавшуюся дырку, поэтому вставка и удаление стоят $O(1)$ без
перевыделения памяти. Тот же приём swap-and-pop применяется и к самим
массивам частиц при удалении --- вот почему любой индекс, переживший
удаление, обязан быть перепроверен.

\subsection{Инкрементальное ведение скоростей}
\label{sec:trick4}

Когда заселённость одного бина меняется на $\delta$,
\[
  T \mathrel{+}= \delta\, \fmaj_{\cdot b},
  \qquad
  \text{row} \gets N \odot (T - \mathrm{diag}\,\fmaj),
  \qquad
  R \gets \textstyle\sum_b \text{row}_b ,
\]
что есть векторная операция $O(B)$ вместо пересчёта полного
матрично-векторного произведения. Выбор использует векторизованные
\code{cumsum} и \code{searchsorted}, а не цикл на Python; иначе именно этот
цикл --- доминирующая стоимость. При очень больших $B$ дерево Фенвика даёт
$O(\log B)$ (Goodson \& Kraft, J.~Comput.\ Phys.\ \textbf{183}, 210 (2002)).

\subsection{Часы}
\label{sec:trick5}

Шаг равен
\begin{equation}
  \Delta t \;=\; \frac{2V}{w\,R},
  \label{eq:num_clock}
\end{equation}
и три отдельные вещи в нём легко испортить.

\paragraph{Упорядоченные и неупорядоченные пары.}
$R$ считает каждую неупорядоченную пару дважды, уравнение~\eqref{eq:num_R},
тогда как физическая мажорантная скорость пробегает неупорядоченные пары.
Отсюда в числителе $2V$, а не $V$. На \emph{выбор} пары это не влияет ---
оба порядка ведут к одной и той же неупорядоченной паре, и множители
сокращаются, --- поэтому ошибка невидима в динамике, перемасштабирует $t$
ровно вдвое и ломает любое сравнение с аналитическим решением.

\paragraph{Веса.}
Здесь $w$ --- константа, но он всё равно вписан в шаг, чтобы формула в
точности совпадала со взвешенным вариантом и чтобы постоянное $w \neq 1$
просто перемасштабировало плотность. Ничто в этом варианте никогда не
меняет $w$, поэтому часы не могут быть испорчены учётом весов. Ради этого
базовый вариант и существует.

\paragraph{Квадратичность.}
$R \propto N^2$, потому что процесс парный, следовательно
\begin{equation}
  t \;\simeq\; \frac{2E}{N^2}
  \label{eq:num_t_of_E}
\end{equation}
после $E$ попыток. Шаг, линейный по $N$, описывает унарный процесс и даёт
неверное $N(t)$. У уравнения~\eqref{eq:num_t_of_E} есть практическое
следствие, на которое легко наткнуться случайно: \emph{увеличение $N$ при
фиксированном бюджете событий делает прогон короче по физическому времени,
а не длиннее}. Любые ворота, заданные абсолютным временем и настроенные при
одном $N$, при $10N$ не пересекаются никогда.

\paragraph{Время идёт на каждой попытке, а не на каждом событии.}
\code{t\_phys} увеличивается сразу после вычисления $\Delta t$, за полсотни
строк до проверки на принятие, так что нулевые (прореженные) попытки тоже
двигают часы. Именно это заставляет работать сокращение
\eqref{eq:num_cancel}, и перенос этого увеличения внутрь ветки принятия ---
та единственная правка, которая превращает мажоранту из стоимости в
\emph{смещение}. Физическое время продвигалось бы тогда на $2V/(wR)$ за одно
принятое событие, поэтому темп событий вышел бы равным $R$ вместо
$R\,\langle K/\fmaj\rangle = \sum K$, и при глобальном $f_{\max}$
\[
  \Delta t_{\rm эфф} \;\propto\; \frac{1}{N^2 f_{\max}}
                     \;\propto\; \frac{1}{N^2 m_{\max}^{\lambda}} .
\]
Ядро вошло бы в часы через \emph{экстремум} распределения, а не через
типичную массу. В открытой системе $m_{\max}$ прижат к стоку и не растёт,
так что зависимость от $m_0$ ушла бы из часов целиком, дрейф выродился бы в
$dm_0/dt \propto m_0$ (экспоненциальный, $\beta = 1$), и условие постоянного
потока вернуло бы $\alpha = -2$ при любом ядре. Это третий, чисто численный
путь к тому же отпечатку $-2$, который разд.~\ref{sec:trick8} производит
физически.

\subsection{Детерминированная инжекция по остатку}
\label{sec:trick6}

Инжекция использует дробный накопитель,
\[
  \nu \mathrel{+}= q\,\Delta t/w,
  \qquad k = \lfloor \nu \rfloor,
  \qquad \nu \mathrel{-}= k,
\]
а не пуассоновский розыгрыш. Это убирает дробовой шум на масштабе инжекции.
Это модельный выбор, и он записывается в метаданные вывода, а не остаётся
неявным.

\subsection{Учёт изохрон}
\label{sec:trick7}

Каждая частица несёт время своего рождения; её возраст есть
$\tau = t - t_{\rm born}$. Снапшоты накапливают двумерную гистограмму по
$(\tau, m)$, строки которой суть изохроны, а сумма по строкам ---
усреднённый по времени спектр.

Для \textbf{фрагментации} возраст осколка однозначен: осколки происходят
ровно от одного родителя. Для \textbf{коагуляции} у слившейся частицы два
предка, и правило наследования приходится выбирать. Оно вынесено в
параметр \code{age\_rule} именно потому, что анализ изохрон от него
зависит, а сделанный выбор обязан быть заявлен и проверен:
\[
  t_{\rm born} \;=\;
  \begin{cases}
    \min(t_1,t_2)                    & \text{\code{'min'}: побеждает старший предок,}\\[2pt]
    t_1 \ \text{если}\ m_1 \ge m_2   & \text{\code{'heavier'},}\\[2pt]
    \dfrac{m_1 t_1 + m_2 t_2}{m_1+m_2} & \text{\code{'mass\_weighted'}.}
  \end{cases}
\]

\subsection{Локальность правила фрагментации --- ловушка физики}
\label{sec:trick8}

Если правило звучит как <<ломается более тяжёлая из пары, чем бы в неё ни
попали>>, то скорость разрушения тела массы $m_0$ на степенном фоне
$F(m) \propto m^{\alpha}$ равна
\begin{equation}
  \nu(m_0) \;\sim\; \int_{\msink}^{m_0} F(m')\,K(m_0,m')\,dm'
           \;\sim\; \int_{\msink/m_0} x^{\alpha}\,dx ,
  \qquad x = m'/m_0 ,
\end{equation}
и этот интеграл \emph{расходится на нижнем пределе всякий раз, когда
$\alpha < -1$}. Дрейф тогда задаётся самыми мелкими частицами в ящике ---
масштабом стока, а не $m_0$. Среднеполевое замыкание теряет силу, дрейф
вырождается в $dm_0/d\tau \sim -C m_0$ (то есть $\beta = 1$,
экспоненциальный спад), и измеряемый индекс запирается на $-2$ независимо
от ядра: мнимая универсальность, которая на деле есть граничный эффект.

Лекарство --- то же, что заложил Dohnanyi (1969): порог катастрофического
разрушения. Пылинка не разбивает валун. Требование
\begin{equation}
  m_{\rm small} \;\ge\; f\,m_{\rm large},
  \qquad f = O(0{,}1\text{--}1),
\end{equation}
ограничивает интеграл областью $x \in [f,1]$, которая свободна от масштаба,
восстанавливает локальность и возвращает $\alpha = -(3+\lambda)/2$.

\emph{Замкнутая} система к этому невосприимчива: её автомодельный пакет
узок, так что все партнёры уже порядка $m_0$ и правило локально по
построению. Именно поэтому замкнутые прогоны согласуются с теорией при
$f=0$, а открытые --- нет.

\subsection{Отсутствие ресэмплинга --- и его цена}
\label{sec:num_cost}

Что покупается:
\begin{itemize}
  \item $\sum w$ и $\sum wm$ сохраняются с машинной точностью; выводимый
        \code{mass\_drift} тождественно нулевой, и любое ненулевое значение ---
        настоящая ошибка;
  \item никаких тонкостей с часами вообще, $\Delta t = 2V/R$ при $w=1$;
  \item ничто не может сместить результат, потому что ничто не
        пересэмплируется.
\end{itemize}

Что платится, и это связывающее ограничение для каждого прогона:
\begin{itemize}
  \item \textbf{Коагуляция.} В замкнутом ящике $N = N_i m_i/m_0$, поэтому
        чтобы наверху оставалось $\sim100$ живых частиц, требуется
        $m_0 \lesssim N_i/100$: при $N_i = 10^6$ это четыре декады, при
        $N_i = 10^8$ --- шесть (и $\sim3$~ГБ массивов). Если продавить
        дальше, прогон заканчивается горсткой частиц; при
        $N_i = 2\times10^5$ один тест финишировал с единственной живой
        частицей, а полезное плато сжалось с $5{,}6$ декад до $2{,}7$.
  \item \textbf{Фрагментация.} Хуже, и в другую сторону: $N$ растёт как
        $m_i/m_0$, так что завершённый каскад держит $N_0\,\minj/\msink$
        частиц. Три декады умножают популяцию в тысячу раз, и \emph{прогон
        заканчивает память, а не физика}.
\end{itemize}

Пользуйтесь этим вариантом, когда приоритет --- безупречная опорная точка.
Пользуйтесь взвешенным вариантом, добавляющим размножение и прореживание
популяции, когда приоритет --- диапазон: десять декад там стоили
$4{,}5\times10^5$ событий и дали $\alpha = -1{,}005$ против теоретического
$-1{,}000$. Оба согласуются всюду, где могут работать оба, --- ради этого
они и держатся вместе.

\subsection{Размещение снапшотов}
\label{sec:trick10}

Наблюдаемая в замкнутой системе --- это сумма Римана
$F(m) = \sum_k (dN/dm)(m,t_k)\,\Delta t_k$, и у выборки по $t$ есть две
раздельные работы, которые нельзя смешивать. \textbf{Размещение} решает,
что будет \emph{разрешено}: при данном $m$ по\-дын\-те\-граль\-ное выражение
имеет
пик при $m_0 \sim m$, поэтому каждой декаде $m$ нужны снапшоты в
соответствующей декаде $m_0$. \textbf{Вес} решает, что будет
\emph{вычислено}: он обязан быть фактическим $\Delta t$ между соседними
снапшотами, каким бы ни было размещение.

\begin{center}
\footnotesize
\begin{tabular}{@{}l p{6.6cm} l@{}}
\toprule
размещение & что происходит & измеренное $\alpha$ \\
 & & (теория $-1$) \\
\midrule
\code{events} & число событий на единицу $m_0$ падает как $m_0^{-2}$,
  поэтому $\Delta t$ взрывается; $94\%$ всего веса приходится на один
  снапшот & $-1{,}4 \ldots -2{,}6$ \\
\code{t\_phys}, узкий диапазон & хорошо работает на $\lesssim 2$ декадах
  & $-0{,}97 \ldots -1{,}02$ \\
\code{t\_phys}, широкий & $t \propto m_0^{1-\lambda}$, поэтому равномерность
  по $t$ есть почти равномерность по $m_0$; на десяти декадах нижние пять
  не получают ни одной выборки & $-0{,}04 \pm 0{,}16$ \\
\code{log\_m0} & $k$ выборок на декаду $m_0$, везде
  & $\mathbf{-1{,}005}$, $R^2 = 0{,}9995$ \\
\bottomrule
\end{tabular}
\end{center}

Провал режима \code{events} заслуживает отдельного упоминания, потому что
он бесшумен: при таком размещении искатель плато из
разд.~\ref{sec:num_measure} всё равно сообщает о чистом плато шириной в
декаду со значением $-1{,}375 \pm 0{,}14$. Внутренне согласованное,
полностью неверное измерение. \emph{Подгонка по плато не спасает сломанный
оценщик}, потому что она проверяет, есть ли у вычисленной вами кривой
скейлинговая область, а не является ли эта кривая правильной величиной.
Независимая гарантия --- печатаемый вес последнего снапшота: он обязан быть
$\ll 1\%$.

Ресэмплинг (разд.~\ref{sec:num_cost}) и размещение лечат разные болезни.
Ресэмплинг покупает \emph{диапазон}; размещение покупает \emph{корректное
измерение} на этом диапазоне.

\paragraph{Правило.} Замкнутая система $\Rightarrow$ \code{log\_m0}, шаг $=$
число снапшотов на декаду ($20$--$30$). Открытая система $\Rightarrow$
стационарное состояние есть \emph{последний} снапшот, $\Delta t$ в оценщик
не входит вовсе, и годится любое расписание.

\section{Сток как критерий}
\label{sec:num_sink}

Критерий остановки или открытия ворот, сформулированный в событиях, есть
утверждение о \emph{ресурсах}; сформулированный в поглощениях на стоке ---
утверждение о \emph{физике}. Только второй переживает смену $N$, диапазона
масс или ядра. Поэтому движок считает $n_{\rm out}$ --- число физических
частиц, покинувших систему через поглощающую границу, --- и выставляет его
наружу тремя способами.

\begin{center}
\small
\begin{tabular}{@{}l p{8.2cm}@{}}
\toprule
\code{iso\_start\_sink} (по умолчанию $10$) &
  начинать накопление изохрон только после $n_{\rm out}$ поглощений \\
\code{stop\_sink\_events} (по умолчанию \code{None}) &
  останавливать прогон на $n_{\rm out}$ поглощениях \\
\code{out["n\_out"]} &
  сам счётчик как временной ряд по снапшотам \\
\bottomrule
\end{tabular}
\end{center}

Ворота говорят: не биновать ни одного возраста, пока каскад доказуемо не
протолкнул столько-то частиц через \emph{весь} инерционный интервал. Их
стоимость следует из баланса массы --- ящик держит $\simeq N\sqrt{\msink}$ в
стационаре, а $n_{\rm out}$ поглощений уносят $n_{\rm out}\msink$, --- так
что
\begin{equation}
  E \;\simeq\; N\sqrt{\msink} \;+\; n_{\rm out}\,\msink .
  \label{eq:num_budget}
\end{equation}
Второе слагаемое не зависит от $N$: ворота не дорожают с ростом прогона,
дорожает только заполнение ящика. Уравнение~\eqref{eq:num_budget} выведено
при $\lambda = 0$ и переоценивает стоимость при $\lambda > 0$, где более
крутое ядро гонит массу к стоку быстрее в расчёте на событие.

У замкнутой системы стока нет, поэтому взведённое значение по умолчанию там
молча разоружается; исключение бросается только на явно переданное
значение. Различие между <<это набрал пользователь>> и <<это пришло из
сигнатуры>> несёт подкласс \code{int}, так что во всех прочих местах
значение ведёт себя как обычное целое.

\section{Измерительный слой}
\label{sec:num_measure}

Прогон производит массивы; превращение их в показатель степени --- отдельная
стадия со своими режимами отказа.

\paragraph{Оценщик.} Замкнутая система: про\-ин\-те\-гри\-ро\-ван\-ная по
возрасту суперпозиция $F(m) = \sum_t (dN/dm)\,\Delta t$. Открытая система:
стационарное состояние \emph{и есть} мгновенный спектр, поэтому оценщиком
служит последний снапшот, а $\Delta t$ в него не входит.

\paragraph{Локальный наклон и плато.} Честный способ смотреть на измеренный
спектр --- это
\begin{equation}
  \Gamma(m) \;=\; \frac{d\log F}{d\log m},
\end{equation}
вычисленный методом наименьших квадратов на скользящем окне. Настоящий
инерционный интервал --- это \emph{плато} в $\Gamma$; загрязнённые концы
проявляются как места, где $\Gamma$ загибается. Сообщаемый индекс есть
среднее $\Gamma$ по самому длинному непрерывному участку, на котором оно
остаётся плоским в пределах допуска, вместе с его разбросом и шириной в
декадах. Плато уже примерно одной декады не является степенным законом,
сколь бы узкой ни была его погрешность; ширина печатается рядом с индексом
именно поэтому. Если ни один участок не проходит, ответ --- \code{nan}, и
это правильный ответ, а не отказ.

\paragraph{Защитная полоса (guard band).} Независимо от этого: отрезать по
$0{,}5$~декады с каждого конца интервала между двумя характерными массами
задачи --- $[\minj,\msink]$ для открытой системы, $[m_i, m_0^{\max}]$ для
замкнутой --- и подогнать там единственный степенной закон. Эта полоса
выбирается \emph{до} того, как посмотрели на данные. Две оценки обязаны
согласоваться; если нет, каскад не развился на том диапазоне, который был
предположен.

Подгонка по всему массиву не даёт ни того, ни другого. На одном прогоне
замкнутой коагуляции при $\lambda=0$ и теоретическом $\alpha=-1$: весь
массив $-1{,}71$, защитная полоса $-1{,}09$, плато $-1{,}16 \pm 0{,}15$ на
$1{,}3$ декады.

\paragraph{Компенсированное представление.} На панелях рисуется
\begin{equation}
  m^2 \frac{dN}{dm} \;=\; \text{масса на логарифмический интервал массы},
\end{equation}
чей наклон равен $\alpha+2$. Вся подгонка делается по сырому $dN/dm$;
компенсация применяется только при отрисовке. Модельная линия несёт индекс
плато и амплитуду, полученную как свободный член наименьших квадратов при
\emph{фиксированном} наклоне, так что никакая вторая подгонка не может
тихо сдвинуть показатель, а теоретическая линия привязывается тем же
способом --- тогда две кривые различаются только наклоном.

\paragraph{Предсказание.} При степени однородности ядра $\lambda$
\begin{equation}
  \beta =
    \begin{cases}
      \lambda & \text{замкнутая: пакет } n \propto m_0^{-2},\\[2pt]
      \dfrac{1+\lambda}{2} & \text{открытая: фон } n \propto m_0^{\alpha},
    \end{cases}
  \qquad
  \alpha = -(1+\beta),
  \quad
  b = \frac{1}{1-\beta}.
\end{equation}
Коагуляция и фрагментация разделяют $\beta$, $b$ и $\alpha$ в точности;
переворачивается только знак дрейфа, поэтому закон читается в $\tau$ при
движении вверх и в $\tau_\ast - \tau$ при движении вниз.

\begin{center}
\footnotesize
\begin{tabular}{@{}@{}l@{\ \ }l@{\ \ }c@{\ \ }c@{\ \ }c@{}}
\toprule
ядро & $K(m_1,m_2)$ & $\lambda$ &
  замк.\ $\alpha$, $b$ &
  откр.\ $\alpha$, $b$ \\
\midrule
постоянное     & $1$                                & $0$   & $-1$,\ \ $1$    & $-3/2$,\ \ $2$ \\
сумма          & $m_1^{1/3}+m_2^{1/3}$              & $1/3$ & $-4/3$,\ \ $3/2$& $-5/3$,\ \ $3$ \\
геометрич.\    & $(m_1^{1/3}+m_2^{1/3})^2$          & $2/3$ & $-5/3$,\ \ $3$  & $-11/6$,\ \ $6$ \\
аддитивное     & $m_1+m_2$                          & $1$   & $-2$,\ \ $\infty$ & $-2$,\ \ $\infty$ \\
произвед.\     & $m_1 m_2$                          & $2$   & \multicolumn{2}{c}{выше гелеобразования} \\
\bottomrule
\end{tabular}
\end{center}

\section{Валидация}
\label{sec:num_validation}

Постоянное ядро, $\lambda = 0$. \emph{плато} --- автоматический индекс,
\emph{полоса} --- априорная защитная полоса, \emph{дек} --- ширина плато в
декадах.

\begin{center}
\small
\begin{tabular}{@{}l r r r r r r@{}}
\toprule
случай & $b$ & $b_{\rm теор}$ & плато & $\pm$ & дек & полоса \\
\midrule
1 замкнутая коагуляция     & $0{,}990$ & $1$ & $-1{,}077$ & $0{,}126$ & $2{,}70$ & $-1{,}035$ \\
2 замкнутая фрагментация   & $1{,}001$ & $1$ & $-0{,}978$ & $0{,}141$ & $3{,}00$ & $-1{,}056$ \\
3 открытая коагуляция      & $2{,}001$ & $2$ & $-1{,}475$ & $0{,}095$ & $3{,}60$ & $-1{,}498$ \\
4 открытая фрагм., $f=0$   & ---       & $2$ & $-2{,}047$ & $0{,}183$ & $0{,}80$ & $-2{,}354$ \\
4 открытая фрагм., $f=0{,}3$ & ---     & $2$ & $-1{,}550$ & $0{,}091$ & $3{,}40$ & $-1{,}517$ \\
\bottomrule
\end{tabular}
\end{center}

Случай~1 дополнительно воспроизводит точное решение Смолуховского для
постоянного ядра, $m_0(t) = m_i(1+n_0K_1t/2)$, с точностью
$\max|{\rm отношение}-1| = 7{,}8\times10^{-14}$ --- прямая проверка
уравнения~\eqref{eq:num_clock}. Дрейф массы тождественно нулевой во всех
четырёх случаях.

Две строки следует читать как диагноз, а не как измерение. Случай~4 при
$f=0$ возвращает $-2{,}05$ на $0{,}80$ декады: индекс залип на $-2$ в
точности так, как предсказывает разд.~\ref{sec:trick8}, а узкое плато само
объявляет, что это не измерение. И показатель роста $b$ для случая~4
отсутствует потому, что при таком бюджете более $99\%$ живых частиц всё ещё
происходят от начального условия и наследуют $t_{\rm born}=0$; ось возраста
тогда вырождена --- возраст равен астрономическому времени сразу у всех, ---
и изохроны деградируют в <<спектр как функция времени>>. Тот же факт
проявляется как $M_{\rm out}/M_{\rm in} > 1$: сток сливает начальный
резервуар, чья масса в $M_{\rm in}$ никогда не входила.

\section{Практические режимы отказа}
\label{sec:num_pitfalls}

\paragraph{Масса вне сетки --- жёсткое падение.}
Мажоранта строится по \emph{верхним углам} бинов, а $b(m)$ зажимает всё,
что выше верхней границы, в последний бин, чей угол тогда меньше самой
частицы. Следовательно $K(m,m) > \fmaj_{B-1,B-1}$, и первая же пара бросает
\code{Majorant violated}. При постоянном ядре этого не может случиться в
принципе, поскольку $K\equiv1$ плоское; при любом $\lambda>0$ это
происходит немедленно. Проверяйте, что сетка охватывает $\minj$ и $\msink$,
до запуска, а не после.

\paragraph{Неограниченный бюджет событий безопасен ровно в одном случае.}
Снимать ограничение на число событий можно только тогда, когда прогон
способно завершить что-то другое:

\begin{center}
\small
\begin{tabular}{@{}l l p{7.4cm}@{}}
\toprule
случай & без предела & что останавливает --- или не останавливает \\
\midrule
замкнутая коагуляция   & безопасно & \code{stop\_max\_mass}; к тому же $N$ монотонно
                                     убывает, так что \code{live}~$<2$ --- гарантированный
                                     терминатор \\
замкнутая фрагментация & опасно    & $N$ растёт; сторож зависания срабатывает лишь
                                     тогда, когда \emph{всё} окажется ниже пассивного
                                     пола, то есть при $N_0 m_i/m_{\rm floor}$ частицах \\
открытая коагуляция    & зависнет  & $N$ сидит на полке, поэтому никакое другое
                                     условие сработать не может \\
открытая фрагментация  & опасно    & $N$ растёт неограниченно \\
\bottomrule
\end{tabular}
\end{center}

В открытых случаях правильный тормоз --- \code{stop\_sink\_events}
(разд.~\ref{sec:num_sink}); дополнительно оставленный ресурсный потолок не
стоит ничего и ограничивает последствия неверно настроенной инжекции.
Записанная причина остановки затем говорит, какой из двух сработал.

\paragraph{Смена ядра обесценивает четыре константы.}
Однородность $\lambda$ и все предсказанные показатели подхватываются
автоматически, а искатель плато, защитная полоса и компенсация вообще не
знают про ядро. Четыре вещи --- знают:

\begin{enumerate}
  \item \textbf{Темп инжекции.} Баланс числа в стационаре есть
        $q = \langle K\rangle N^2/(2V)$. Запись $q = N^2/2$ --- это та же
        формула при $\langle K\rangle = 1$. Поскольку $N$ набирается на
        масштабе инжекции при $\alpha<-1$, большинство пар суть
        $(\minj,\minj)$, так что $\langle K \rangle \sim K(\minj,\minj)$, а
        более тяжёлые партнёры поднимают эту величину примерно вдвое.
  \item \textbf{Время столкновения.} $t_c = 1/(Kn)$, а не $1/n$.
  \item \textbf{Биновка по возрасту.} Один возрастной бин шириной
        $\Delta_\tau$ декад покрывает $b\,\Delta_\tau$ декад \emph{массы}.
        Фиксируйте разрешение по массе и выводите шаг по возрасту из него,
        $\Delta_\tau = 0{,}30/b$; иначе при $b=6$ шаг, настроенный под
        $b=2$, перепрыгивает почти декаду массы за бин, и двухдекадное окно
        подгонки ловит два бина.
  \item \textbf{Точные решения.} Замкнутая форма Смолуховского и проверка
        <<$m_0(t)$ обязано быть линейным>> --- свойства исключительно
        постоянного ядра.
\end{enumerate}

\paragraph{У отпечатка $-2$ три разные причины.}
Индекс $-2$ возвращают: (i) настоящее ядро с $\lambda = 1$, где это
правильный ответ; (ii) нелокальное правило разрушения из
разд.~\ref{sec:trick8}, где это граничный эффект; (iii) часы, продвигаемые
на событие, а не на попытку, разд.~\ref{sec:trick5}, где это численный
артефакт. Все три гонят $\beta \to 1$ и на бумаге выглядят одинаково. Само
по себе значение $-2$ поэтому не говорит ни о чём. Различающая проверка ---
меняется ли индекс при смене ядра: если не меняется, это не универсальность.

\paragraph{Плоская полка --- не инерционный интервал.}
В замкнутой фрагментации пассивный пол замораживает всё, что ниже него: эти
частицы больше никогда не дробятся и накапливаются в плоскую полку, к
каскаду отношения не имеющую. Оставленная в массиве, эта полка --- самый
длинный плоский участок из всех, $4{,}3$ декады при $\alpha \simeq 0$, --- и
искатель плато послушно сообщает именно о ней. Ограничивайте поиск областью
$m \ge m_{\rm floor}$.

% ---------------------------------------- удалить при вставке в main.tex
\end{document}

%      рядом с уже имеющимся \include{Parts/unified-cascade_v4}.
%      В преамбулу main.tex тогда нужно добавить:
%          \usepackage{algorithm}
%          \usepackage{algpseudocode}
%          \usepackage{booktabs}
%      и, поскольку main.tex сейчас английский, переключить babel на
%          \usepackage[english,russian]{babel}
%      и кодировку на \usepackage[T1,T2A]{fontenc}.
% =====================================================================

% ------------------------------------------ ПРЕАМБУЛА -- ОТРЕЗАТЬ ОТСЮДА
\documentclass[a4paper,12pt]{article}

\usepackage[T1,T2A]{fontenc}
\usepackage[utf8]{inputenc}
\usepackage[english,russian]{babel}
\usepackage{amsmath,amssymb}
\usepackage{geometry}
\geometry{margin=2.5cm}
\usepackage{graphicx}
\usepackage{booktabs}
\usepackage{algorithm}
\usepackage{algpseudocode}
\usepackage[unicode,colorlinks=true,linkcolor=blue,citecolor=blue,urlcolor=blue]{hyperref}

\newcommand{\minj}{m_{\mathrm{inj}}}
\newcommand{\msink}{m_{\mathrm{sink}}}
\newcommand{\fmaj}{f^{\mathrm{maj}}}
\newcommand{\code}[1]{\texttt{#1}}
\emergencystretch=2em   % русские слова длиннее английских

\begin{document}
\selectlanguage{russian}

\title{Численная схема: метод Монте-Карло с мажорантной частотой\\ для каскада в пространстве масс}
\author{}
\date{}
\maketitle
% ------------------------------------------ ПРЕАМБУЛА -- ОТРЕЗАТЬ ДОСЮДА

\section{Область применимости и принцип построения}
\label{sec:num_scope}

Движок интегрирует стохастический каскад в пространстве масс в четырёх
конфигурациях,
\[
  \text{процесс} \in \{\text{коагуляция},\ \text{фрагментация}\},
  \qquad
  \text{система} \in \{\text{замкнутая},\ \text{открытая}\},
\]
одним и тем же путём в коде. Между четырьмя случаями различаются ровно две
вещи:

\begin{enumerate}
  \item что принятое событие делает с двумя выбранными частицами, и
  \item граничные условия --- инжекция плюс поглощающий сток, либо ни того,
        ни другого.
\end{enumerate}

Всё остальное --- построение мажоранты, выбор пары, часы, учёт ---
общее. Это сознательное ограничение, а не удобство: расхождение между
двумя из четырёх случаев тогда нельзя списать на два разных интегратора.

Описываемый здесь вариант не несёт весов частиц. Одна расчётная частица
есть одна физическая частица, $w \equiv 1$ на весь прогон, поэтому
$\sum_i w_i$ и $\sum_i w_i m_i$ сохраняются с машинной точностью на каждом
шаге. Измеряемый дрейф массы тождественно равен нулю, и любое ненулевое
значение --- это ошибка в коде, а не статистическое отклонение.
Раздел~\ref{sec:num_cost} говорит, чего это стоит.

\subsection{Состояние}

Состояние --- четыре параллельных массива по живым частицам плюс разбиение
пространства масс на бины:

\begin{center}
\small
\begin{tabular}{@{}ll@{}}
\toprule
\code{mass[i]}       & масса частицы $m_i$ \\
\code{inj\_time[i]}  & время рождения, нужно для изохрон (разд.~\ref{sec:trick7}) \\
\code{bin\_of[i]}    & номер массового бина, в котором лежит $i$ \\
\code{pos\_in\_bin[i]} & позиция $i$ внутри списка индексов этого бина \\
\midrule
\code{edges[0..B]}   & $B+1$ логарифмически расположенных границ бинов \\
\code{bins[b]}       & список индексов частиц в бине $b$ \\
\code{N\_bin[b]}     & заселённость $N_b$ \\
\bottomrule
\end{tabular}
\end{center}

Частица массы $m$ попадает в бин
$b(m) = \max\bigl(0,\min(B-1,\ \mathrm{searchsorted}(\text{edges},m)-1)\bigr)$.
Зажим (clamping) в этом выражении --- ловушка, и она разобрана в
разд.~\ref{sec:num_pitfalls}.

\section{Базовый алгоритм}
\label{sec:num_algorithm}

\subsection{Инициализация}

\begin{algorithm}[h]
\caption{Инициализация: бины, таблица мажоранты, вектор скоростей}
\label{alg:setup}
\begin{algorithmic}[1]
\State разложить все начальные частицы по бинам; сформировать заселённости $N_b$
\State \textbf{таблица мажоранты:} для всех $b,c$ \quad
       $\fmaj_{bc} \gets K(\text{edges}[b+1],\ \text{edges}[c+1])$
       \Comment{верхние углы}
\State $T_b \gets \sum_c \fmaj_{bc} N_c$
\State $\text{row}_b \gets N_b\,(T_b - \fmaj_{bb})$
\State $R \gets \sum_b \text{row}_b$
\end{algorithmic}
\end{algorithm}

Угловое правило $\fmaj_{bc}=K(\text{edges}[b+1],\text{edges}[c+1])$
ограничивает $K$ сверху на всей паре бинов \emph{при условии, что $K$ не
убывает по обоим аргументам}. Это единственное структурное требование к
ядру, и оно проверяется во время счёта, а не предполагается: каждая
принятая попытка проверяет $K(m_i,m_j) \le \fmaj_{bc}$ и падает, если это
нарушено.

По построению
\begin{equation}
  R \;=\; \sum_b N_b\bigl(T_b - \fmaj_{bb}\bigr)
    \;=\; \sum_{i \neq j} \fmaj_{ij}
    \;=\; 2 \sum_{i<j} \fmaj_{ij},
  \label{eq:num_R}
\end{equation}
то есть $R$ пробегает \emph{упорядоченные} пары. Двойка в
уравнении~\eqref{eq:num_R} --- источник двойки в часах,
уравнение~\eqref{eq:num_clock}.

\subsection{Главный цикл}

\begin{algorithm}[h]
\caption{Главный цикл: одна мажорантная попытка}
\label{alg:mainloop}
\begin{algorithmic}[1]
\While{ни одно условие остановки не выполнено}
  \State $\Delta t \gets 2V/(wR)$; \quad $t \gets t + \Delta t$
         \Comment{часы, разд.~\ref{sec:trick5}}
  \State \textbf{инжекция (открытые системы):}
         $\nu \gets \nu + q\,\Delta t / w$; \quad
         $k \gets \lfloor \nu \rfloor$; \quad $\nu \gets \nu - k$;
         добавить $k$ частиц массы $\minj$
  \State выбрать бин-строку $b$ с $P(b) \propto \text{row}_b$
  \State выбрать бин-столбец $c$ с
         $P(c) \propto N_c \fmaj_{bc}$, беря $(N_b-1)\fmaj_{bb}$ при $c=b$
  \State выбрать $i$ равномерно из \code{bins[b]}, $j \neq i$ равномерно из \code{bins[c]}
  \State \textbf{прореживание:} принять с вероятностью $K(m_i,m_j)/\fmaj_{bc}$;
         иначе \textbf{continue} \Comment{нулевое событие}
  \State применить событие (алг.~\ref{alg:event}); применить поглощающий сток
  \State записать снапшот, если этого требует расписание (разд.~\ref{sec:trick10})
\EndWhile
\end{algorithmic}
\end{algorithm}

Отклонённая попытка --- не ошибка: именно прореживание делает доминирующий
процесс точным. Доля принятия, выводимая как \code{acc} в логе прогона,
есть $O(1)$ по построению бин-парной таблицы --- измерено $0{,}95$ для
геометрического ядра на бинах шириной $0{,}1$~декады, где угловой перебор
составляет $10^{0.1\lambda}$ на аргумент.

\subsection{Событие}

\begin{algorithm}[h]
\caption{Принятое событие --- единственное место, где четыре случая расходятся}
\label{alg:event}
\begin{algorithmic}[1]
\If{коагуляция}
  \State $m_{\mathrm{new}} \gets m_i + m_j$
  \State унаследовать возраст $i$ от пары $(i,j)$ по \code{age\_rule} (разд.~\ref{sec:trick7})
  \State $m_i \gets m_{\mathrm{new}}$; перебиновать $i$; \textbf{удалить} $j$
  \State \textbf{если} открытая \textbf{и} $m_{\mathrm{new}} \ge \msink$ \textbf{то}
         удалить $i$; $M_{\mathrm{out}} \mathrel{+}= w\,m_{\mathrm{new}}$;
         $n_{\mathrm{out}} \mathrel{+}= 1$
\Else \Comment{фрагментация}
  \State $p \gets \arg\max(m_i,m_j)$; \quad $s \gets \arg\min(m_i,m_j)$
  \State \textbf{если} $m_s < f\,m_p$ \textbf{то continue}
         \Comment{порог разрушения, разд.~\ref{sec:trick8}}
  \State $\xi \sim U(0,1)$; \quad $m_a \gets \xi m_p$, \ $m_b \gets (1-\xi)m_p$
  \State оба осколка наследуют $t_{\mathrm{born}}$ единственного родителя $p$
  \State $m_p \gets m_a$; перебиновать $p$; \textbf{добавить} частицу массы $m_b$
  \State \textbf{если} открытая: удалить каждый осколок с $m \le \msink$,
         начиная со старшего индекса
\EndIf
\end{algorithmic}
\end{algorithm}

Обратите внимание на асимметрию знака --- ею управляется всё содержимое
разд.~\ref{sec:num_pitfalls}: коагуляция есть $2 \to 1$ и, следовательно,
\emph{сток} числа частиц, фрагментация есть $1 \to 2$ и, следовательно,
\emph{источник}.

\section{Десять трюков}
\label{sec:num_tricks}

\subsection{Мажорантная частота столкновений}
\label{sec:trick1}

Вычислять точную парную скорость для каждой пары --- это $O(N^2)$ на шаг.
Вместо этого мажоранта $\fmaj \ge K$ управляет доминирующим пуассоновским
процессом, а события прореживаются с вероятностью принятия $K/\fmaj$.
Прореженный процесс точен: мажоранта влияет на стоимость, но никогда --- на
статистику.

\emph{Источники.} Ivanov \& Rogasinsky, Sov.\ J.\ Numer.\ Anal.\ Math.\
Modelling \textbf{3}, 453 (1988) [схема MCF в DSMC]; Garcia, van den
Broeck, Aertsens \& Serneels, Physica~A \textbf{143}, 535 (1987)
[мажоранта для коагуляции]; Gillespie, J.~Atmos.\ Sci.\ \textbf{32}, 1977
(1975) [точная стохастическая коалесценция].

\subsection{Бин-парная таблица мажоранты}
\label{sec:trick2}

Мажоранта --- это матрица $\fmaj_{bc}$ по парам массовых бинов, а не одно
глобальное число. Покупается этим \emph{стоимость} и \emph{проверяемость}.
Стоит проговорить, что именно, потому что корректность мажорантой купить
нельзя: она сокращается из физического темпа в точности.

\paragraph{Мажоранта сокращается.}
Тики доминирующего процесса приходят с темпом $wR/(2V)$; тик выбирает
упорядоченную пару $(i,j)$ с вероятностью $\fmaj_{ij}/R$; принимается он с
вероятностью $K_{ij}/\fmaj_{ij}$. Темп принятых событий на этой паре, стало
быть, равен
\begin{equation}
  \frac{wR}{2V}\cdot\frac{\fmaj_{ij}}{R}\cdot\frac{K_{ij}}{\fmaj_{ij}}
  \;=\; \frac{w\,K_{ij}}{2V},
  \label{eq:num_cancel}
\end{equation}
а суммирование по неупорядоченным парам даёт $w\sum_{i<j}K_{ij}/V$ ---
смолуховский темп, вообще не зависящий от $\fmaj$. Мажоранта сокращается
дважды: один раз в весе выбора пары, второй раз в доле принятия.
Следовательно, \emph{любая} допустимая мажоранта возвращает одну и ту же
физику; плохая лишь впустую тратит попытки. Равенство~\eqref{eq:num_cancel}
верно потому, что часы продвигаются на каждой \emph{попытке}, а не на каждом
событии --- см.\ разд.~\ref{sec:trick5}, где разобрана та единственная
правка, которая это ломает.

\paragraph{Что таблица покупает на самом деле.}
Одно глобальное $f_{\max} = K(m_{\max},m_{\max})$ --- вполне законная
мажоранта, дающая ровно ту же статистику. И тем не менее это неверный выбор,
по трём причинам:

\begin{itemize}
  \item \textbf{Доля принятия.} Типичная пара принимается с вероятностью
        $\langle K\rangle/f_{\max} \simeq (m_0/m_{\max})^{\lambda}$. На $D$
        декадах это $10^{-\lambda D}$: для геометрического ядра на четырёх
        декадах --- $2\times10^{-3}$, то есть около пятисот попыток на
        событие против $1{,}05$ с таблицей. Таблица же держит долю принятия
        $O(1)$ в \emph{каждой} паре бинов, потому что граница берётся в углах
        именно той пары, которую разыграли; перебор тогда составляет всего
        $10^{0.1\lambda}$ на аргумент при бинах в $0{,}1$~декады.
  \item \textbf{Ложное срабатывание сторожа зависания.}
        \code{max\_stall\_tries} считает попытки без принятого события. При
        доле принятия $2\times10^{-3}$ его умолчание $2\times10^{6}$ ---
        это всего $\sim\!4000$ событий честной работы, так что здоровый
        прогон может быть убит с диагнозом <<acceptance collapsed>>.
  \item \textbf{Отстающий максимум --- настоящее смещение.} \emph{Бегущий}
        глобальный максимум, не обновлённый до того, как какая-то частица
        переросла его, даёт $\fmaj < K$, и прореживание тогда молча неверно;
        это, в отличие от двух пунктов выше, уже отказ по корректности.
        Угловая таблица статична, и каждая попытка явно проверяет
        $K \le \fmaj$, поэтому та же ситуация приводит к падению, а не к
        смещению.
\end{itemize}

\emph{Источник.} Eibeck \& Wagner, SIAM J.~Sci.\ Comput.\ \textbf{22}, 802
(2000),
\url{https://doi.org/10.1137/S1064827599353488}.

\subsection{Членство в бине за $O(1)$}
\label{sec:trick3}

\code{bins[b]} --- список индексов частиц, а \code{pos\_in\_bin[i]} ---
позиция $i$ внутри него. Удаление переставляет последний элемент в
образовавшуюся дырку, поэтому вставка и удаление стоят $O(1)$ без
перевыделения памяти. Тот же приём swap-and-pop применяется и к самим
массивам частиц при удалении --- вот почему любой индекс, переживший
удаление, обязан быть перепроверен.

\subsection{Инкрементальное ведение скоростей}
\label{sec:trick4}

Когда заселённость одного бина меняется на $\delta$,
\[
  T \mathrel{+}= \delta\, \fmaj_{\cdot b},
  \qquad
  \text{row} \gets N \odot (T - \mathrm{diag}\,\fmaj),
  \qquad
  R \gets \textstyle\sum_b \text{row}_b ,
\]
что есть векторная операция $O(B)$ вместо пересчёта полного
матрично-векторного произведения. Выбор использует векторизованные
\code{cumsum} и \code{searchsorted}, а не цикл на Python; иначе именно этот
цикл --- доминирующая стоимость. При очень больших $B$ дерево Фенвика даёт
$O(\log B)$ (Goodson \& Kraft, J.~Comput.\ Phys.\ \textbf{183}, 210 (2002)).

\subsection{Часы}
\label{sec:trick5}

Шаг равен
\begin{equation}
  \Delta t \;=\; \frac{2V}{w\,R},
  \label{eq:num_clock}
\end{equation}
и три отдельные вещи в нём легко испортить.

\paragraph{Упорядоченные и неупорядоченные пары.}
$R$ считает каждую неупорядоченную пару дважды, уравнение~\eqref{eq:num_R},
тогда как физическая мажорантная скорость пробегает неупорядоченные пары.
Отсюда в числителе $2V$, а не $V$. На \emph{выбор} пары это не влияет ---
оба порядка ведут к одной и той же неупорядоченной паре, и множители
сокращаются, --- поэтому ошибка невидима в динамике, перемасштабирует $t$
ровно вдвое и ломает любое сравнение с аналитическим решением.

\paragraph{Веса.}
Здесь $w$ --- константа, но он всё равно вписан в шаг, чтобы формула в
точности совпадала со взвешенным вариантом и чтобы постоянное $w \neq 1$
просто перемасштабировало плотность. Ничто в этом варианте никогда не
меняет $w$, поэтому часы не могут быть испорчены учётом весов. Ради этого
базовый вариант и существует.

\paragraph{Квадратичность.}
$R \propto N^2$, потому что процесс парный, следовательно
\begin{equation}
  t \;\simeq\; \frac{2E}{N^2}
  \label{eq:num_t_of_E}
\end{equation}
после $E$ попыток. Шаг, линейный по $N$, описывает унарный процесс и даёт
неверное $N(t)$. У уравнения~\eqref{eq:num_t_of_E} есть практическое
следствие, на которое легко наткнуться случайно: \emph{увеличение $N$ при
фиксированном бюджете событий делает прогон короче по физическому времени,
а не длиннее}. Любые ворота, заданные абсолютным временем и настроенные при
одном $N$, при $10N$ не пересекаются никогда.

\paragraph{Время идёт на каждой попытке, а не на каждом событии.}
\code{t\_phys} увеличивается сразу после вычисления $\Delta t$, за полсотни
строк до проверки на принятие, так что нулевые (прореженные) попытки тоже
двигают часы. Именно это заставляет работать сокращение
\eqref{eq:num_cancel}, и перенос этого увеличения внутрь ветки принятия ---
та единственная правка, которая превращает мажоранту из стоимости в
\emph{смещение}. Физическое время продвигалось бы тогда на $2V/(wR)$ за одно
принятое событие, поэтому темп событий вышел бы равным $R$ вместо
$R\,\langle K/\fmaj\rangle = \sum K$, и при глобальном $f_{\max}$
\[
  \Delta t_{\rm эфф} \;\propto\; \frac{1}{N^2 f_{\max}}
                     \;\propto\; \frac{1}{N^2 m_{\max}^{\lambda}} .
\]
Ядро вошло бы в часы через \emph{экстремум} распределения, а не через
типичную массу. В открытой системе $m_{\max}$ прижат к стоку и не растёт,
так что зависимость от $m_0$ ушла бы из часов целиком, дрейф выродился бы в
$dm_0/dt \propto m_0$ (экспоненциальный, $\beta = 1$), и условие постоянного
потока вернуло бы $\alpha = -2$ при любом ядре. Это третий, чисто численный
путь к тому же отпечатку $-2$, который разд.~\ref{sec:trick8} производит
физически.

\subsection{Детерминированная инжекция по остатку}
\label{sec:trick6}

Инжекция использует дробный накопитель,
\[
  \nu \mathrel{+}= q\,\Delta t/w,
  \qquad k = \lfloor \nu \rfloor,
  \qquad \nu \mathrel{-}= k,
\]
а не пуассоновский розыгрыш. Это убирает дробовой шум на масштабе инжекции.
Это модельный выбор, и он записывается в метаданные вывода, а не остаётся
неявным.

\subsection{Учёт изохрон}
\label{sec:trick7}

Каждая частица несёт время своего рождения; её возраст есть
$\tau = t - t_{\rm born}$. Снапшоты накапливают двумерную гистограмму по
$(\tau, m)$, строки которой суть изохроны, а сумма по строкам ---
усреднённый по времени спектр.

Для \textbf{фрагментации} возраст осколка однозначен: осколки происходят
ровно от одного родителя. Для \textbf{коагуляции} у слившейся частицы два
предка, и правило наследования приходится выбирать. Оно вынесено в
параметр \code{age\_rule} именно потому, что анализ изохрон от него
зависит, а сделанный выбор обязан быть заявлен и проверен:
\[
  t_{\rm born} \;=\;
  \begin{cases}
    \min(t_1,t_2)                    & \text{\code{'min'}: побеждает старший предок,}\\[2pt]
    t_1 \ \text{если}\ m_1 \ge m_2   & \text{\code{'heavier'},}\\[2pt]
    \dfrac{m_1 t_1 + m_2 t_2}{m_1+m_2} & \text{\code{'mass\_weighted'}.}
  \end{cases}
\]

\subsection{Локальность правила фрагментации --- ловушка физики}
\label{sec:trick8}

Если правило звучит как <<ломается более тяжёлая из пары, чем бы в неё ни
попали>>, то скорость разрушения тела массы $m_0$ на степенном фоне
$F(m) \propto m^{\alpha}$ равна
\begin{equation}
  \nu(m_0) \;\sim\; \int_{\msink}^{m_0} F(m')\,K(m_0,m')\,dm'
           \;\sim\; \int_{\msink/m_0} x^{\alpha}\,dx ,
  \qquad x = m'/m_0 ,
\end{equation}
и этот интеграл \emph{расходится на нижнем пределе всякий раз, когда
$\alpha < -1$}. Дрейф тогда задаётся самыми мелкими частицами в ящике ---
масштабом стока, а не $m_0$. Среднеполевое замыкание теряет силу, дрейф
вырождается в $dm_0/d\tau \sim -C m_0$ (то есть $\beta = 1$,
экспоненциальный спад), и измеряемый индекс запирается на $-2$ независимо
от ядра: мнимая универсальность, которая на деле есть граничный эффект.

Лекарство --- то же, что заложил Dohnanyi (1969): порог катастрофического
разрушения. Пылинка не разбивает валун. Требование
\begin{equation}
  m_{\rm small} \;\ge\; f\,m_{\rm large},
  \qquad f = O(0{,}1\text{--}1),
\end{equation}
ограничивает интеграл областью $x \in [f,1]$, которая свободна от масштаба,
восстанавливает локальность и возвращает $\alpha = -(3+\lambda)/2$.

\emph{Замкнутая} система к этому невосприимчива: её автомодельный пакет
узок, так что все партнёры уже порядка $m_0$ и правило локально по
построению. Именно поэтому замкнутые прогоны согласуются с теорией при
$f=0$, а открытые --- нет.

\subsection{Отсутствие ресэмплинга --- и его цена}
\label{sec:num_cost}

Что покупается:
\begin{itemize}
  \item $\sum w$ и $\sum wm$ сохраняются с машинной точностью; выводимый
        \code{mass\_drift} тождественно нулевой, и любое ненулевое значение ---
        настоящая ошибка;
  \item никаких тонкостей с часами вообще, $\Delta t = 2V/R$ при $w=1$;
  \item ничто не может сместить результат, потому что ничто не
        пересэмплируется.
\end{itemize}

Что платится, и это связывающее ограничение для каждого прогона:
\begin{itemize}
  \item \textbf{Коагуляция.} В замкнутом ящике $N = N_i m_i/m_0$, поэтому
        чтобы наверху оставалось $\sim100$ живых частиц, требуется
        $m_0 \lesssim N_i/100$: при $N_i = 10^6$ это четыре декады, при
        $N_i = 10^8$ --- шесть (и $\sim3$~ГБ массивов). Если продавить
        дальше, прогон заканчивается горсткой частиц; при
        $N_i = 2\times10^5$ один тест финишировал с единственной живой
        частицей, а полезное плато сжалось с $5{,}6$ декад до $2{,}7$.
  \item \textbf{Фрагментация.} Хуже, и в другую сторону: $N$ растёт как
        $m_i/m_0$, так что завершённый каскад держит $N_0\,\minj/\msink$
        частиц. Три декады умножают популяцию в тысячу раз, и \emph{прогон
        заканчивает память, а не физика}.
\end{itemize}

Пользуйтесь этим вариантом, когда приоритет --- безупречная опорная точка.
Пользуйтесь взвешенным вариантом, добавляющим размножение и прореживание
популяции, когда приоритет --- диапазон: десять декад там стоили
$4{,}5\times10^5$ событий и дали $\alpha = -1{,}005$ против теоретического
$-1{,}000$. Оба согласуются всюду, где могут работать оба, --- ради этого
они и держатся вместе.

\subsection{Размещение снапшотов}
\label{sec:trick10}

Наблюдаемая в замкнутой системе --- это сумма Римана
$F(m) = \sum_k (dN/dm)(m,t_k)\,\Delta t_k$, и у выборки по $t$ есть две
раздельные работы, которые нельзя смешивать. \textbf{Размещение} решает,
что будет \emph{разрешено}: при данном $m$ по\-дын\-те\-граль\-ное выражение
имеет
пик при $m_0 \sim m$, поэтому каждой декаде $m$ нужны снапшоты в
соответствующей декаде $m_0$. \textbf{Вес} решает, что будет
\emph{вычислено}: он обязан быть фактическим $\Delta t$ между соседними
снапшотами, каким бы ни было размещение.

\begin{center}
\footnotesize
\begin{tabular}{@{}l p{6.6cm} l@{}}
\toprule
размещение & что происходит & измеренное $\alpha$ \\
 & & (теория $-1$) \\
\midrule
\code{events} & число событий на единицу $m_0$ падает как $m_0^{-2}$,
  поэтому $\Delta t$ взрывается; $94\%$ всего веса приходится на один
  снапшот & $-1{,}4 \ldots -2{,}6$ \\
\code{t\_phys}, узкий диапазон & хорошо работает на $\lesssim 2$ декадах
  & $-0{,}97 \ldots -1{,}02$ \\
\code{t\_phys}, широкий & $t \propto m_0^{1-\lambda}$, поэтому равномерность
  по $t$ есть почти равномерность по $m_0$; на десяти декадах нижние пять
  не получают ни одной выборки & $-0{,}04 \pm 0{,}16$ \\
\code{log\_m0} & $k$ выборок на декаду $m_0$, везде
  & $\mathbf{-1{,}005}$, $R^2 = 0{,}9995$ \\
\bottomrule
\end{tabular}
\end{center}

Провал режима \code{events} заслуживает отдельного упоминания, потому что
он бесшумен: при таком размещении искатель плато из
разд.~\ref{sec:num_measure} всё равно сообщает о чистом плато шириной в
декаду со значением $-1{,}375 \pm 0{,}14$. Внутренне согласованное,
полностью неверное измерение. \emph{Подгонка по плато не спасает сломанный
оценщик}, потому что она проверяет, есть ли у вычисленной вами кривой
скейлинговая область, а не является ли эта кривая правильной величиной.
Независимая гарантия --- печатаемый вес последнего снапшота: он обязан быть
$\ll 1\%$.

Ресэмплинг (разд.~\ref{sec:num_cost}) и размещение лечат разные болезни.
Ресэмплинг покупает \emph{диапазон}; размещение покупает \emph{корректное
измерение} на этом диапазоне.

\paragraph{Правило.} Замкнутая система $\Rightarrow$ \code{log\_m0}, шаг $=$
число снапшотов на декаду ($20$--$30$). Открытая система $\Rightarrow$
стационарное состояние есть \emph{последний} снапшот, $\Delta t$ в оценщик
не входит вовсе, и годится любое расписание.

\section{Сток как критерий}
\label{sec:num_sink}

Критерий остановки или открытия ворот, сформулированный в событиях, есть
утверждение о \emph{ресурсах}; сформулированный в поглощениях на стоке ---
утверждение о \emph{физике}. Только второй переживает смену $N$, диапазона
масс или ядра. Поэтому движок считает $n_{\rm out}$ --- число физических
частиц, покинувших систему через поглощающую границу, --- и выставляет его
наружу тремя способами.

\begin{center}
\small
\begin{tabular}{@{}l p{8.2cm}@{}}
\toprule
\code{iso\_start\_sink} (по умолчанию $10$) &
  начинать накопление изохрон только после $n_{\rm out}$ поглощений \\
\code{stop\_sink\_events} (по умолчанию \code{None}) &
  останавливать прогон на $n_{\rm out}$ поглощениях \\
\code{out["n\_out"]} &
  сам счётчик как временной ряд по снапшотам \\
\bottomrule
\end{tabular}
\end{center}

Ворота говорят: не биновать ни одного возраста, пока каскад доказуемо не
протолкнул столько-то частиц через \emph{весь} инерционный интервал. Их
стоимость следует из баланса массы --- ящик держит $\simeq N\sqrt{\msink}$ в
стационаре, а $n_{\rm out}$ поглощений уносят $n_{\rm out}\msink$, --- так
что
\begin{equation}
  E \;\simeq\; N\sqrt{\msink} \;+\; n_{\rm out}\,\msink .
  \label{eq:num_budget}
\end{equation}
Второе слагаемое не зависит от $N$: ворота не дорожают с ростом прогона,
дорожает только заполнение ящика. Уравнение~\eqref{eq:num_budget} выведено
при $\lambda = 0$ и переоценивает стоимость при $\lambda > 0$, где более
крутое ядро гонит массу к стоку быстрее в расчёте на событие.

У замкнутой системы стока нет, поэтому взведённое значение по умолчанию там
молча разоружается; исключение бросается только на явно переданное
значение. Различие между <<это набрал пользователь>> и <<это пришло из
сигнатуры>> несёт подкласс \code{int}, так что во всех прочих местах
значение ведёт себя как обычное целое.

\section{Измерительный слой}
\label{sec:num_measure}

Прогон производит массивы; превращение их в показатель степени --- отдельная
стадия со своими режимами отказа.

\paragraph{Оценщик.} Замкнутая система: про\-ин\-те\-гри\-ро\-ван\-ная по
возрасту суперпозиция $F(m) = \sum_t (dN/dm)\,\Delta t$. Открытая система:
стационарное состояние \emph{и есть} мгновенный спектр, поэтому оценщиком
служит последний снапшот, а $\Delta t$ в него не входит.

\paragraph{Локальный наклон и плато.} Честный способ смотреть на измеренный
спектр --- это
\begin{equation}
  \Gamma(m) \;=\; \frac{d\log F}{d\log m},
\end{equation}
вычисленный методом наименьших квадратов на скользящем окне. Настоящий
инерционный интервал --- это \emph{плато} в $\Gamma$; загрязнённые концы
проявляются как места, где $\Gamma$ загибается. Сообщаемый индекс есть
среднее $\Gamma$ по самому длинному непрерывному участку, на котором оно
остаётся плоским в пределах допуска, вместе с его разбросом и шириной в
декадах. Плато уже примерно одной декады не является степенным законом,
сколь бы узкой ни была его погрешность; ширина печатается рядом с индексом
именно поэтому. Если ни один участок не проходит, ответ --- \code{nan}, и
это правильный ответ, а не отказ.

\paragraph{Защитная полоса (guard band).} Независимо от этого: отрезать по
$0{,}5$~декады с каждого конца интервала между двумя характерными массами
задачи --- $[\minj,\msink]$ для открытой системы, $[m_i, m_0^{\max}]$ для
замкнутой --- и подогнать там единственный степенной закон. Эта полоса
выбирается \emph{до} того, как посмотрели на данные. Две оценки обязаны
согласоваться; если нет, каскад не развился на том диапазоне, который был
предположен.

Подгонка по всему массиву не даёт ни того, ни другого. На одном прогоне
замкнутой коагуляции при $\lambda=0$ и теоретическом $\alpha=-1$: весь
массив $-1{,}71$, защитная полоса $-1{,}09$, плато $-1{,}16 \pm 0{,}15$ на
$1{,}3$ декады.

\paragraph{Компенсированное представление.} На панелях рисуется
\begin{equation}
  m^2 \frac{dN}{dm} \;=\; \text{масса на логарифмический интервал массы},
\end{equation}
чей наклон равен $\alpha+2$. Вся подгонка делается по сырому $dN/dm$;
компенсация применяется только при отрисовке. Модельная линия несёт индекс
плато и амплитуду, полученную как свободный член наименьших квадратов при
\emph{фиксированном} наклоне, так что никакая вторая подгонка не может
тихо сдвинуть показатель, а теоретическая линия привязывается тем же
способом --- тогда две кривые различаются только наклоном.

\paragraph{Предсказание.} При степени однородности ядра $\lambda$
\begin{equation}
  \beta =
    \begin{cases}
      \lambda & \text{замкнутая: пакет } n \propto m_0^{-2},\\[2pt]
      \dfrac{1+\lambda}{2} & \text{открытая: фон } n \propto m_0^{\alpha},
    \end{cases}
  \qquad
  \alpha = -(1+\beta),
  \quad
  b = \frac{1}{1-\beta}.
\end{equation}
Коагуляция и фрагментация разделяют $\beta$, $b$ и $\alpha$ в точности;
переворачивается только знак дрейфа, поэтому закон читается в $\tau$ при
движении вверх и в $\tau_\ast - \tau$ при движении вниз.

\begin{center}
\footnotesize
\begin{tabular}{@{}@{}l@{\ \ }l@{\ \ }c@{\ \ }c@{\ \ }c@{}}
\toprule
ядро & $K(m_1,m_2)$ & $\lambda$ &
  замк.\ $\alpha$, $b$ &
  откр.\ $\alpha$, $b$ \\
\midrule
постоянное     & $1$                                & $0$   & $-1$,\ \ $1$    & $-3/2$,\ \ $2$ \\
сумма          & $m_1^{1/3}+m_2^{1/3}$              & $1/3$ & $-4/3$,\ \ $3/2$& $-5/3$,\ \ $3$ \\
геометрич.\    & $(m_1^{1/3}+m_2^{1/3})^2$          & $2/3$ & $-5/3$,\ \ $3$  & $-11/6$,\ \ $6$ \\
аддитивное     & $m_1+m_2$                          & $1$   & $-2$,\ \ $\infty$ & $-2$,\ \ $\infty$ \\
произвед.\     & $m_1 m_2$                          & $2$   & \multicolumn{2}{c}{выше гелеобразования} \\
\bottomrule
\end{tabular}
\end{center}

\section{Валидация}
\label{sec:num_validation}

Постоянное ядро, $\lambda = 0$. \emph{плато} --- автоматический индекс,
\emph{полоса} --- априорная защитная полоса, \emph{дек} --- ширина плато в
декадах.

\begin{center}
\small
\begin{tabular}{@{}l r r r r r r@{}}
\toprule
случай & $b$ & $b_{\rm теор}$ & плато & $\pm$ & дек & полоса \\
\midrule
1 замкнутая коагуляция     & $0{,}990$ & $1$ & $-1{,}077$ & $0{,}126$ & $2{,}70$ & $-1{,}035$ \\
2 замкнутая фрагментация   & $1{,}001$ & $1$ & $-0{,}978$ & $0{,}141$ & $3{,}00$ & $-1{,}056$ \\
3 открытая коагуляция      & $2{,}001$ & $2$ & $-1{,}475$ & $0{,}095$ & $3{,}60$ & $-1{,}498$ \\
4 открытая фрагм., $f=0$   & ---       & $2$ & $-2{,}047$ & $0{,}183$ & $0{,}80$ & $-2{,}354$ \\
4 открытая фрагм., $f=0{,}3$ & ---     & $2$ & $-1{,}550$ & $0{,}091$ & $3{,}40$ & $-1{,}517$ \\
\bottomrule
\end{tabular}
\end{center}

Случай~1 дополнительно воспроизводит точное решение Смолуховского для
постоянного ядра, $m_0(t) = m_i(1+n_0K_1t/2)$, с точностью
$\max|{\rm отношение}-1| = 7{,}8\times10^{-14}$ --- прямая проверка
уравнения~\eqref{eq:num_clock}. Дрейф массы тождественно нулевой во всех
четырёх случаях.

Две строки следует читать как диагноз, а не как измерение. Случай~4 при
$f=0$ возвращает $-2{,}05$ на $0{,}80$ декады: индекс залип на $-2$ в
точности так, как предсказывает разд.~\ref{sec:trick8}, а узкое плато само
объявляет, что это не измерение. И показатель роста $b$ для случая~4
отсутствует потому, что при таком бюджете более $99\%$ живых частиц всё ещё
происходят от начального условия и наследуют $t_{\rm born}=0$; ось возраста
тогда вырождена --- возраст равен астрономическому времени сразу у всех, ---
и изохроны деградируют в <<спектр как функция времени>>. Тот же факт
проявляется как $M_{\rm out}/M_{\rm in} > 1$: сток сливает начальный
резервуар, чья масса в $M_{\rm in}$ никогда не входила.

\section{Практические режимы отказа}
\label{sec:num_pitfalls}

\paragraph{Масса вне сетки --- жёсткое падение.}
Мажоранта строится по \emph{верхним углам} бинов, а $b(m)$ зажимает всё,
что выше верхней границы, в последний бин, чей угол тогда меньше самой
частицы. Следовательно $K(m,m) > \fmaj_{B-1,B-1}$, и первая же пара бросает
\code{Majorant violated}. При постоянном ядре этого не может случиться в
принципе, поскольку $K\equiv1$ плоское; при любом $\lambda>0$ это
происходит немедленно. Проверяйте, что сетка охватывает $\minj$ и $\msink$,
до запуска, а не после.

\paragraph{Неограниченный бюджет событий безопасен ровно в одном случае.}
Снимать ограничение на число событий можно только тогда, когда прогон
способно завершить что-то другое:

\begin{center}
\small
\begin{tabular}{@{}l l p{7.4cm}@{}}
\toprule
случай & без предела & что останавливает --- или не останавливает \\
\midrule
замкнутая коагуляция   & безопасно & \code{stop\_max\_mass}; к тому же $N$ монотонно
                                     убывает, так что \code{live}~$<2$ --- гарантированный
                                     терминатор \\
замкнутая фрагментация & опасно    & $N$ растёт; сторож зависания срабатывает лишь
                                     тогда, когда \emph{всё} окажется ниже пассивного
                                     пола, то есть при $N_0 m_i/m_{\rm floor}$ частицах \\
открытая коагуляция    & зависнет  & $N$ сидит на полке, поэтому никакое другое
                                     условие сработать не может \\
открытая фрагментация  & опасно    & $N$ растёт неограниченно \\
\bottomrule
\end{tabular}
\end{center}

В открытых случаях правильный тормоз --- \code{stop\_sink\_events}
(разд.~\ref{sec:num_sink}); дополнительно оставленный ресурсный потолок не
стоит ничего и ограничивает последствия неверно настроенной инжекции.
Записанная причина остановки затем говорит, какой из двух сработал.

\paragraph{Смена ядра обесценивает четыре константы.}
Однородность $\lambda$ и все предсказанные показатели подхватываются
автоматически, а искатель плато, защитная полоса и компенсация вообще не
знают про ядро. Четыре вещи --- знают:

\begin{enumerate}
  \item \textbf{Темп инжекции.} Баланс числа в стационаре есть
        $q = \langle K\rangle N^2/(2V)$. Запись $q = N^2/2$ --- это та же
        формула при $\langle K\rangle = 1$. Поскольку $N$ набирается на
        масштабе инжекции при $\alpha<-1$, большинство пар суть
        $(\minj,\minj)$, так что $\langle K \rangle \sim K(\minj,\minj)$, а
        более тяжёлые партнёры поднимают эту величину примерно вдвое.
  \item \textbf{Время столкновения.} $t_c = 1/(Kn)$, а не $1/n$.
  \item \textbf{Биновка по возрасту.} Один возрастной бин шириной
        $\Delta_\tau$ декад покрывает $b\,\Delta_\tau$ декад \emph{массы}.
        Фиксируйте разрешение по массе и выводите шаг по возрасту из него,
        $\Delta_\tau = 0{,}30/b$; иначе при $b=6$ шаг, настроенный под
        $b=2$, перепрыгивает почти декаду массы за бин, и двухдекадное окно
        подгонки ловит два бина.
  \item \textbf{Точные решения.} Замкнутая форма Смолуховского и проверка
        <<$m_0(t)$ обязано быть линейным>> --- свойства исключительно
        постоянного ядра.
\end{enumerate}

\paragraph{У отпечатка $-2$ три разные причины.}
Индекс $-2$ возвращают: (i) настоящее ядро с $\lambda = 1$, где это
правильный ответ; (ii) нелокальное правило разрушения из
разд.~\ref{sec:trick8}, где это граничный эффект; (iii) часы, продвигаемые
на событие, а не на попытку, разд.~\ref{sec:trick5}, где это численный
артефакт. Все три гонят $\beta \to 1$ и на бумаге выглядят одинаково. Само
по себе значение $-2$ поэтому не говорит ни о чём. Различающая проверка ---
меняется ли индекс при смене ядра: если не меняется, это не универсальность.

\paragraph{Плоская полка --- не инерционный интервал.}
В замкнутой фрагментации пассивный пол замораживает всё, что ниже него: эти
частицы больше никогда не дробятся и накапливаются в плоскую полку, к
каскаду отношения не имеющую. Оставленная в массиве, эта полка --- самый
длинный плоский участок из всех, $4{,}3$ декады при $\alpha \simeq 0$, --- и
искатель плато послушно сообщает именно о ней. Ограничивайте поиск областью
$m \ge m_{\rm floor}$.

% ---------------------------------------- удалить при вставке в main.tex
\end{document}

%      рядом с уже имеющимся \include{Parts/unified-cascade_v4}.
%      В преамбулу main.tex тогда нужно добавить:
%          \usepackage{algorithm}
%          \usepackage{algpseudocode}
%          \usepackage{booktabs}
%      и, поскольку main.tex сейчас английский, переключить babel на
%          \usepackage[english,russian]{babel}
%      и кодировку на \usepackage[T1,T2A]{fontenc}.
% =====================================================================

% ------------------------------------------ ПРЕАМБУЛА -- ОТРЕЗАТЬ ОТСЮДА
\documentclass[a4paper,12pt]{article}

\usepackage[T1,T2A]{fontenc}
\usepackage[utf8]{inputenc}
\usepackage[english,russian]{babel}
\usepackage{amsmath,amssymb}
\usepackage{geometry}
\geometry{margin=2.5cm}
\usepackage{graphicx}
\usepackage{booktabs}
\usepackage{algorithm}
\usepackage{algpseudocode}
\usepackage[unicode,colorlinks=true,linkcolor=blue,citecolor=blue,urlcolor=blue]{hyperref}

\newcommand{\minj}{m_{\mathrm{inj}}}
\newcommand{\msink}{m_{\mathrm{sink}}}
\newcommand{\fmaj}{f^{\mathrm{maj}}}
\newcommand{\code}[1]{\texttt{#1}}
\emergencystretch=2em   % русские слова длиннее английских

\begin{document}
\selectlanguage{russian}

\title{Численная схема: метод Монте-Карло с мажорантной частотой\\ для каскада в пространстве масс}
\author{}
\date{}
\maketitle
% ------------------------------------------ ПРЕАМБУЛА -- ОТРЕЗАТЬ ДОСЮДА

\section{Область применимости и принцип построения}
\label{sec:num_scope}

Движок интегрирует стохастический каскад в пространстве масс в четырёх
конфигурациях,
\[
  \text{процесс} \in \{\text{коагуляция},\ \text{фрагментация}\},
  \qquad
  \text{система} \in \{\text{замкнутая},\ \text{открытая}\},
\]
одним и тем же путём в коде. Между четырьмя случаями различаются ровно две
вещи:

\begin{enumerate}
  \item что принятое событие делает с двумя выбранными частицами, и
  \item граничные условия --- инжекция плюс поглощающий сток, либо ни того,
        ни другого.
\end{enumerate}

Всё остальное --- построение мажоранты, выбор пары, часы, учёт ---
общее. Это сознательное ограничение, а не удобство: расхождение между
двумя из четырёх случаев тогда нельзя списать на два разных интегратора.

Описываемый здесь вариант не несёт весов частиц. Одна расчётная частица
есть одна физическая частица, $w \equiv 1$ на весь прогон, поэтому
$\sum_i w_i$ и $\sum_i w_i m_i$ сохраняются с машинной точностью на каждом
шаге. Измеряемый дрейф массы тождественно равен нулю, и любое ненулевое
значение --- это ошибка в коде, а не статистическое отклонение.
Раздел~\ref{sec:num_cost} говорит, чего это стоит.

\subsection{Состояние}

Состояние --- четыре параллельных массива по живым частицам плюс разбиение
пространства масс на бины:

\begin{center}
\small
\begin{tabular}{@{}ll@{}}
\toprule
\code{mass[i]}       & масса частицы $m_i$ \\
\code{inj\_time[i]}  & время рождения, нужно для изохрон (разд.~\ref{sec:trick7}) \\
\code{bin\_of[i]}    & номер массового бина, в котором лежит $i$ \\
\code{pos\_in\_bin[i]} & позиция $i$ внутри списка индексов этого бина \\
\midrule
\code{edges[0..B]}   & $B+1$ логарифмически расположенных границ бинов \\
\code{bins[b]}       & список индексов частиц в бине $b$ \\
\code{N\_bin[b]}     & заселённость $N_b$ \\
\bottomrule
\end{tabular}
\end{center}

Частица массы $m$ попадает в бин
$b(m) = \max\bigl(0,\min(B-1,\ \mathrm{searchsorted}(\text{edges},m)-1)\bigr)$.
Зажим (clamping) в этом выражении --- ловушка, и она разобрана в
разд.~\ref{sec:num_pitfalls}.

\section{Базовый алгоритм}
\label{sec:num_algorithm}

\subsection{Инициализация}

\begin{algorithm}[h]
\caption{Инициализация: бины, таблица мажоранты, вектор скоростей}
\label{alg:setup}
\begin{algorithmic}[1]
\State разложить все начальные частицы по бинам; сформировать заселённости $N_b$
\State \textbf{таблица мажоранты:} для всех $b,c$ \quad
       $\fmaj_{bc} \gets K(\text{edges}[b+1],\ \text{edges}[c+1])$
       \Comment{верхние углы}
\State $T_b \gets \sum_c \fmaj_{bc} N_c$
\State $\text{row}_b \gets N_b\,(T_b - \fmaj_{bb})$
\State $R \gets \sum_b \text{row}_b$
\end{algorithmic}
\end{algorithm}

Угловое правило $\fmaj_{bc}=K(\text{edges}[b+1],\text{edges}[c+1])$
ограничивает $K$ сверху на всей паре бинов \emph{при условии, что $K$ не
убывает по обоим аргументам}. Это единственное структурное требование к
ядру, и оно проверяется во время счёта, а не предполагается: каждая
принятая попытка проверяет $K(m_i,m_j) \le \fmaj_{bc}$ и падает, если это
нарушено.

По построению
\begin{equation}
  R \;=\; \sum_b N_b\bigl(T_b - \fmaj_{bb}\bigr)
    \;=\; \sum_{i \neq j} \fmaj_{ij}
    \;=\; 2 \sum_{i<j} \fmaj_{ij},
  \label{eq:num_R}
\end{equation}
то есть $R$ пробегает \emph{упорядоченные} пары. Двойка в
уравнении~\eqref{eq:num_R} --- источник двойки в часах,
уравнение~\eqref{eq:num_clock}.

\subsection{Главный цикл}

\begin{algorithm}[h]
\caption{Главный цикл: одна мажорантная попытка}
\label{alg:mainloop}
\begin{algorithmic}[1]
\While{ни одно условие остановки не выполнено}
  \State $\Delta t \gets 2V/(wR)$; \quad $t \gets t + \Delta t$
         \Comment{часы, разд.~\ref{sec:trick5}}
  \State \textbf{инжекция (открытые системы):}
         $\nu \gets \nu + q\,\Delta t / w$; \quad
         $k \gets \lfloor \nu \rfloor$; \quad $\nu \gets \nu - k$;
         добавить $k$ частиц массы $\minj$
  \State выбрать бин-строку $b$ с $P(b) \propto \text{row}_b$
  \State выбрать бин-столбец $c$ с
         $P(c) \propto N_c \fmaj_{bc}$, беря $(N_b-1)\fmaj_{bb}$ при $c=b$
  \State выбрать $i$ равномерно из \code{bins[b]}, $j \neq i$ равномерно из \code{bins[c]}
  \State \textbf{прореживание:} принять с вероятностью $K(m_i,m_j)/\fmaj_{bc}$;
         иначе \textbf{continue} \Comment{нулевое событие}
  \State применить событие (алг.~\ref{alg:event}); применить поглощающий сток
  \State записать снапшот, если этого требует расписание (разд.~\ref{sec:trick10})
\EndWhile
\end{algorithmic}
\end{algorithm}

Отклонённая попытка --- не ошибка: именно прореживание делает доминирующий
процесс точным. Доля принятия, выводимая как \code{acc} в логе прогона,
есть $O(1)$ по построению бин-парной таблицы --- измерено $0{,}95$ для
геометрического ядра на бинах шириной $0{,}1$~декады, где угловой перебор
составляет $10^{0.1\lambda}$ на аргумент.

\subsection{Событие}

\begin{algorithm}[h]
\caption{Принятое событие --- единственное место, где четыре случая расходятся}
\label{alg:event}
\begin{algorithmic}[1]
\If{коагуляция}
  \State $m_{\mathrm{new}} \gets m_i + m_j$
  \State унаследовать возраст $i$ от пары $(i,j)$ по \code{age\_rule} (разд.~\ref{sec:trick7})
  \State $m_i \gets m_{\mathrm{new}}$; перебиновать $i$; \textbf{удалить} $j$
  \State \textbf{если} открытая \textbf{и} $m_{\mathrm{new}} \ge \msink$ \textbf{то}
         удалить $i$; $M_{\mathrm{out}} \mathrel{+}= w\,m_{\mathrm{new}}$;
         $n_{\mathrm{out}} \mathrel{+}= 1$
\Else \Comment{фрагментация}
  \State $p \gets \arg\max(m_i,m_j)$; \quad $s \gets \arg\min(m_i,m_j)$
  \State \textbf{если} $m_s < f\,m_p$ \textbf{то continue}
         \Comment{порог разрушения, разд.~\ref{sec:trick8}}
  \State $\xi \sim U(0,1)$; \quad $m_a \gets \xi m_p$, \ $m_b \gets (1-\xi)m_p$
  \State оба осколка наследуют $t_{\mathrm{born}}$ единственного родителя $p$
  \State $m_p \gets m_a$; перебиновать $p$; \textbf{добавить} частицу массы $m_b$
  \State \textbf{если} открытая: удалить каждый осколок с $m \le \msink$,
         начиная со старшего индекса
\EndIf
\end{algorithmic}
\end{algorithm}

Обратите внимание на асимметрию знака --- ею управляется всё содержимое
разд.~\ref{sec:num_pitfalls}: коагуляция есть $2 \to 1$ и, следовательно,
\emph{сток} числа частиц, фрагментация есть $1 \to 2$ и, следовательно,
\emph{источник}.

\section{Десять трюков}
\label{sec:num_tricks}

\subsection{Мажорантная частота столкновений}
\label{sec:trick1}

Вычислять точную парную скорость для каждой пары --- это $O(N^2)$ на шаг.
Вместо этого мажоранта $\fmaj \ge K$ управляет доминирующим пуассоновским
процессом, а события прореживаются с вероятностью принятия $K/\fmaj$.
Прореженный процесс точен: мажоранта влияет на стоимость, но никогда --- на
статистику.

\emph{Источники.} Ivanov \& Rogasinsky, Sov.\ J.\ Numer.\ Anal.\ Math.\
Modelling \textbf{3}, 453 (1988) [схема MCF в DSMC]; Garcia, van den
Broeck, Aertsens \& Serneels, Physica~A \textbf{143}, 535 (1987)
[мажоранта для коагуляции]; Gillespie, J.~Atmos.\ Sci.\ \textbf{32}, 1977
(1975) [точная стохастическая коалесценция].

\subsection{Бин-парная таблица мажоранты}
\label{sec:trick2}

Мажоранта --- это матрица $\fmaj_{bc}$ по парам массовых бинов, а не одно
глобальное число. Покупается этим \emph{стоимость} и \emph{проверяемость}.
Стоит проговорить, что именно, потому что корректность мажорантой купить
нельзя: она сокращается из физического темпа в точности.

\paragraph{Мажоранта сокращается.}
Тики доминирующего процесса приходят с темпом $wR/(2V)$; тик выбирает
упорядоченную пару $(i,j)$ с вероятностью $\fmaj_{ij}/R$; принимается он с
вероятностью $K_{ij}/\fmaj_{ij}$. Темп принятых событий на этой паре, стало
быть, равен
\begin{equation}
  \frac{wR}{2V}\cdot\frac{\fmaj_{ij}}{R}\cdot\frac{K_{ij}}{\fmaj_{ij}}
  \;=\; \frac{w\,K_{ij}}{2V},
  \label{eq:num_cancel}
\end{equation}
а суммирование по неупорядоченным парам даёт $w\sum_{i<j}K_{ij}/V$ ---
смолуховский темп, вообще не зависящий от $\fmaj$. Мажоранта сокращается
дважды: один раз в весе выбора пары, второй раз в доле принятия.
Следовательно, \emph{любая} допустимая мажоранта возвращает одну и ту же
физику; плохая лишь впустую тратит попытки. Равенство~\eqref{eq:num_cancel}
верно потому, что часы продвигаются на каждой \emph{попытке}, а не на каждом
событии --- см.\ разд.~\ref{sec:trick5}, где разобрана та единственная
правка, которая это ломает.

\paragraph{Что таблица покупает на самом деле.}
Одно глобальное $f_{\max} = K(m_{\max},m_{\max})$ --- вполне законная
мажоранта, дающая ровно ту же статистику. И тем не менее это неверный выбор,
по трём причинам:

\begin{itemize}
  \item \textbf{Доля принятия.} Типичная пара принимается с вероятностью
        $\langle K\rangle/f_{\max} \simeq (m_0/m_{\max})^{\lambda}$. На $D$
        декадах это $10^{-\lambda D}$: для геометрического ядра на четырёх
        декадах --- $2\times10^{-3}$, то есть около пятисот попыток на
        событие против $1{,}05$ с таблицей. Таблица же держит долю принятия
        $O(1)$ в \emph{каждой} паре бинов, потому что граница берётся в углах
        именно той пары, которую разыграли; перебор тогда составляет всего
        $10^{0.1\lambda}$ на аргумент при бинах в $0{,}1$~декады.
  \item \textbf{Ложное срабатывание сторожа зависания.}
        \code{max\_stall\_tries} считает попытки без принятого события. При
        доле принятия $2\times10^{-3}$ его умолчание $2\times10^{6}$ ---
        это всего $\sim\!4000$ событий честной работы, так что здоровый
        прогон может быть убит с диагнозом <<acceptance collapsed>>.
  \item \textbf{Отстающий максимум --- настоящее смещение.} \emph{Бегущий}
        глобальный максимум, не обновлённый до того, как какая-то частица
        переросла его, даёт $\fmaj < K$, и прореживание тогда молча неверно;
        это, в отличие от двух пунктов выше, уже отказ по корректности.
        Угловая таблица статична, и каждая попытка явно проверяет
        $K \le \fmaj$, поэтому та же ситуация приводит к падению, а не к
        смещению.
\end{itemize}

\emph{Источник.} Eibeck \& Wagner, SIAM J.~Sci.\ Comput.\ \textbf{22}, 802
(2000),
\url{https://doi.org/10.1137/S1064827599353488}.

\subsection{Членство в бине за $O(1)$}
\label{sec:trick3}

\code{bins[b]} --- список индексов частиц, а \code{pos\_in\_bin[i]} ---
позиция $i$ внутри него. Удаление переставляет последний элемент в
образовавшуюся дырку, поэтому вставка и удаление стоят $O(1)$ без
перевыделения памяти. Тот же приём swap-and-pop применяется и к самим
массивам частиц при удалении --- вот почему любой индекс, переживший
удаление, обязан быть перепроверен.

\subsection{Инкрементальное ведение скоростей}
\label{sec:trick4}

Когда заселённость одного бина меняется на $\delta$,
\[
  T \mathrel{+}= \delta\, \fmaj_{\cdot b},
  \qquad
  \text{row} \gets N \odot (T - \mathrm{diag}\,\fmaj),
  \qquad
  R \gets \textstyle\sum_b \text{row}_b ,
\]
что есть векторная операция $O(B)$ вместо пересчёта полного
матрично-векторного произведения. Выбор использует векторизованные
\code{cumsum} и \code{searchsorted}, а не цикл на Python; иначе именно этот
цикл --- доминирующая стоимость. При очень больших $B$ дерево Фенвика даёт
$O(\log B)$ (Goodson \& Kraft, J.~Comput.\ Phys.\ \textbf{183}, 210 (2002)).

\subsection{Часы}
\label{sec:trick5}

Шаг равен
\begin{equation}
  \Delta t \;=\; \frac{2V}{w\,R},
  \label{eq:num_clock}
\end{equation}
и три отдельные вещи в нём легко испортить.

\paragraph{Упорядоченные и неупорядоченные пары.}
$R$ считает каждую неупорядоченную пару дважды, уравнение~\eqref{eq:num_R},
тогда как физическая мажорантная скорость пробегает неупорядоченные пары.
Отсюда в числителе $2V$, а не $V$. На \emph{выбор} пары это не влияет ---
оба порядка ведут к одной и той же неупорядоченной паре, и множители
сокращаются, --- поэтому ошибка невидима в динамике, перемасштабирует $t$
ровно вдвое и ломает любое сравнение с аналитическим решением.

\paragraph{Веса.}
Здесь $w$ --- константа, но он всё равно вписан в шаг, чтобы формула в
точности совпадала со взвешенным вариантом и чтобы постоянное $w \neq 1$
просто перемасштабировало плотность. Ничто в этом варианте никогда не
меняет $w$, поэтому часы не могут быть испорчены учётом весов. Ради этого
базовый вариант и существует.

\paragraph{Квадратичность.}
$R \propto N^2$, потому что процесс парный, следовательно
\begin{equation}
  t \;\simeq\; \frac{2E}{N^2}
  \label{eq:num_t_of_E}
\end{equation}
после $E$ попыток. Шаг, линейный по $N$, описывает унарный процесс и даёт
неверное $N(t)$. У уравнения~\eqref{eq:num_t_of_E} есть практическое
следствие, на которое легко наткнуться случайно: \emph{увеличение $N$ при
фиксированном бюджете событий делает прогон короче по физическому времени,
а не длиннее}. Любые ворота, заданные абсолютным временем и настроенные при
одном $N$, при $10N$ не пересекаются никогда.

\paragraph{Время идёт на каждой попытке, а не на каждом событии.}
\code{t\_phys} увеличивается сразу после вычисления $\Delta t$, за полсотни
строк до проверки на принятие, так что нулевые (прореженные) попытки тоже
двигают часы. Именно это заставляет работать сокращение
\eqref{eq:num_cancel}, и перенос этого увеличения внутрь ветки принятия ---
та единственная правка, которая превращает мажоранту из стоимости в
\emph{смещение}. Физическое время продвигалось бы тогда на $2V/(wR)$ за одно
принятое событие, поэтому темп событий вышел бы равным $R$ вместо
$R\,\langle K/\fmaj\rangle = \sum K$, и при глобальном $f_{\max}$
\[
  \Delta t_{\rm эфф} \;\propto\; \frac{1}{N^2 f_{\max}}
                     \;\propto\; \frac{1}{N^2 m_{\max}^{\lambda}} .
\]
Ядро вошло бы в часы через \emph{экстремум} распределения, а не через
типичную массу. В открытой системе $m_{\max}$ прижат к стоку и не растёт,
так что зависимость от $m_0$ ушла бы из часов целиком, дрейф выродился бы в
$dm_0/dt \propto m_0$ (экспоненциальный, $\beta = 1$), и условие постоянного
потока вернуло бы $\alpha = -2$ при любом ядре. Это третий, чисто численный
путь к тому же отпечатку $-2$, который разд.~\ref{sec:trick8} производит
физически.

\subsection{Детерминированная инжекция по остатку}
\label{sec:trick6}

Инжекция использует дробный накопитель,
\[
  \nu \mathrel{+}= q\,\Delta t/w,
  \qquad k = \lfloor \nu \rfloor,
  \qquad \nu \mathrel{-}= k,
\]
а не пуассоновский розыгрыш. Это убирает дробовой шум на масштабе инжекции.
Это модельный выбор, и он записывается в метаданные вывода, а не остаётся
неявным.

\subsection{Учёт изохрон}
\label{sec:trick7}

Каждая частица несёт время своего рождения; её возраст есть
$\tau = t - t_{\rm born}$. Снапшоты накапливают двумерную гистограмму по
$(\tau, m)$, строки которой суть изохроны, а сумма по строкам ---
усреднённый по времени спектр.

Для \textbf{фрагментации} возраст осколка однозначен: осколки происходят
ровно от одного родителя. Для \textbf{коагуляции} у слившейся частицы два
предка, и правило наследования приходится выбирать. Оно вынесено в
параметр \code{age\_rule} именно потому, что анализ изохрон от него
зависит, а сделанный выбор обязан быть заявлен и проверен:
\[
  t_{\rm born} \;=\;
  \begin{cases}
    \min(t_1,t_2)                    & \text{\code{'min'}: побеждает старший предок,}\\[2pt]
    t_1 \ \text{если}\ m_1 \ge m_2   & \text{\code{'heavier'},}\\[2pt]
    \dfrac{m_1 t_1 + m_2 t_2}{m_1+m_2} & \text{\code{'mass\_weighted'}.}
  \end{cases}
\]

\subsection{Локальность правила фрагментации --- ловушка физики}
\label{sec:trick8}

Если правило звучит как <<ломается более тяжёлая из пары, чем бы в неё ни
попали>>, то скорость разрушения тела массы $m_0$ на степенном фоне
$F(m) \propto m^{\alpha}$ равна
\begin{equation}
  \nu(m_0) \;\sim\; \int_{\msink}^{m_0} F(m')\,K(m_0,m')\,dm'
           \;\sim\; \int_{\msink/m_0} x^{\alpha}\,dx ,
  \qquad x = m'/m_0 ,
\end{equation}
и этот интеграл \emph{расходится на нижнем пределе всякий раз, когда
$\alpha < -1$}. Дрейф тогда задаётся самыми мелкими частицами в ящике ---
масштабом стока, а не $m_0$. Среднеполевое замыкание теряет силу, дрейф
вырождается в $dm_0/d\tau \sim -C m_0$ (то есть $\beta = 1$,
экспоненциальный спад), и измеряемый индекс запирается на $-2$ независимо
от ядра: мнимая универсальность, которая на деле есть граничный эффект.

Лекарство --- то же, что заложил Dohnanyi (1969): порог катастрофического
разрушения. Пылинка не разбивает валун. Требование
\begin{equation}
  m_{\rm small} \;\ge\; f\,m_{\rm large},
  \qquad f = O(0{,}1\text{--}1),
\end{equation}
ограничивает интеграл областью $x \in [f,1]$, которая свободна от масштаба,
восстанавливает локальность и возвращает $\alpha = -(3+\lambda)/2$.

\emph{Замкнутая} система к этому невосприимчива: её автомодельный пакет
узок, так что все партнёры уже порядка $m_0$ и правило локально по
построению. Именно поэтому замкнутые прогоны согласуются с теорией при
$f=0$, а открытые --- нет.

\subsection{Отсутствие ресэмплинга --- и его цена}
\label{sec:num_cost}

Что покупается:
\begin{itemize}
  \item $\sum w$ и $\sum wm$ сохраняются с машинной точностью; выводимый
        \code{mass\_drift} тождественно нулевой, и любое ненулевое значение ---
        настоящая ошибка;
  \item никаких тонкостей с часами вообще, $\Delta t = 2V/R$ при $w=1$;
  \item ничто не может сместить результат, потому что ничто не
        пересэмплируется.
\end{itemize}

Что платится, и это связывающее ограничение для каждого прогона:
\begin{itemize}
  \item \textbf{Коагуляция.} В замкнутом ящике $N = N_i m_i/m_0$, поэтому
        чтобы наверху оставалось $\sim100$ живых частиц, требуется
        $m_0 \lesssim N_i/100$: при $N_i = 10^6$ это четыре декады, при
        $N_i = 10^8$ --- шесть (и $\sim3$~ГБ массивов). Если продавить
        дальше, прогон заканчивается горсткой частиц; при
        $N_i = 2\times10^5$ один тест финишировал с единственной живой
        частицей, а полезное плато сжалось с $5{,}6$ декад до $2{,}7$.
  \item \textbf{Фрагментация.} Хуже, и в другую сторону: $N$ растёт как
        $m_i/m_0$, так что завершённый каскад держит $N_0\,\minj/\msink$
        частиц. Три декады умножают популяцию в тысячу раз, и \emph{прогон
        заканчивает память, а не физика}.
\end{itemize}

Пользуйтесь этим вариантом, когда приоритет --- безупречная опорная точка.
Пользуйтесь взвешенным вариантом, добавляющим размножение и прореживание
популяции, когда приоритет --- диапазон: десять декад там стоили
$4{,}5\times10^5$ событий и дали $\alpha = -1{,}005$ против теоретического
$-1{,}000$. Оба согласуются всюду, где могут работать оба, --- ради этого
они и держатся вместе.

\subsection{Размещение снапшотов}
\label{sec:trick10}

Наблюдаемая в замкнутой системе --- это сумма Римана
$F(m) = \sum_k (dN/dm)(m,t_k)\,\Delta t_k$, и у выборки по $t$ есть две
раздельные работы, которые нельзя смешивать. \textbf{Размещение} решает,
что будет \emph{разрешено}: при данном $m$ по\-дын\-те\-граль\-ное выражение
имеет
пик при $m_0 \sim m$, поэтому каждой декаде $m$ нужны снапшоты в
соответствующей декаде $m_0$. \textbf{Вес} решает, что будет
\emph{вычислено}: он обязан быть фактическим $\Delta t$ между соседними
снапшотами, каким бы ни было размещение.

\begin{center}
\footnotesize
\begin{tabular}{@{}l p{6.6cm} l@{}}
\toprule
размещение & что происходит & измеренное $\alpha$ \\
 & & (теория $-1$) \\
\midrule
\code{events} & число событий на единицу $m_0$ падает как $m_0^{-2}$,
  поэтому $\Delta t$ взрывается; $94\%$ всего веса приходится на один
  снапшот & $-1{,}4 \ldots -2{,}6$ \\
\code{t\_phys}, узкий диапазон & хорошо работает на $\lesssim 2$ декадах
  & $-0{,}97 \ldots -1{,}02$ \\
\code{t\_phys}, широкий & $t \propto m_0^{1-\lambda}$, поэтому равномерность
  по $t$ есть почти равномерность по $m_0$; на десяти декадах нижние пять
  не получают ни одной выборки & $-0{,}04 \pm 0{,}16$ \\
\code{log\_m0} & $k$ выборок на декаду $m_0$, везде
  & $\mathbf{-1{,}005}$, $R^2 = 0{,}9995$ \\
\bottomrule
\end{tabular}
\end{center}

Провал режима \code{events} заслуживает отдельного упоминания, потому что
он бесшумен: при таком размещении искатель плато из
разд.~\ref{sec:num_measure} всё равно сообщает о чистом плато шириной в
декаду со значением $-1{,}375 \pm 0{,}14$. Внутренне согласованное,
полностью неверное измерение. \emph{Подгонка по плато не спасает сломанный
оценщик}, потому что она проверяет, есть ли у вычисленной вами кривой
скейлинговая область, а не является ли эта кривая правильной величиной.
Независимая гарантия --- печатаемый вес последнего снапшота: он обязан быть
$\ll 1\%$.

Ресэмплинг (разд.~\ref{sec:num_cost}) и размещение лечат разные болезни.
Ресэмплинг покупает \emph{диапазон}; размещение покупает \emph{корректное
измерение} на этом диапазоне.

\paragraph{Правило.} Замкнутая система $\Rightarrow$ \code{log\_m0}, шаг $=$
число снапшотов на декаду ($20$--$30$). Открытая система $\Rightarrow$
стационарное состояние есть \emph{последний} снапшот, $\Delta t$ в оценщик
не входит вовсе, и годится любое расписание.

\section{Сток как критерий}
\label{sec:num_sink}

Критерий остановки или открытия ворот, сформулированный в событиях, есть
утверждение о \emph{ресурсах}; сформулированный в поглощениях на стоке ---
утверждение о \emph{физике}. Только второй переживает смену $N$, диапазона
масс или ядра. Поэтому движок считает $n_{\rm out}$ --- число физических
частиц, покинувших систему через поглощающую границу, --- и выставляет его
наружу тремя способами.

\begin{center}
\small
\begin{tabular}{@{}l p{8.2cm}@{}}
\toprule
\code{iso\_start\_sink} (по умолчанию $10$) &
  начинать накопление изохрон только после $n_{\rm out}$ поглощений \\
\code{stop\_sink\_events} (по умолчанию \code{None}) &
  останавливать прогон на $n_{\rm out}$ поглощениях \\
\code{out["n\_out"]} &
  сам счётчик как временной ряд по снапшотам \\
\bottomrule
\end{tabular}
\end{center}

Ворота говорят: не биновать ни одного возраста, пока каскад доказуемо не
протолкнул столько-то частиц через \emph{весь} инерционный интервал. Их
стоимость следует из баланса массы --- ящик держит $\simeq N\sqrt{\msink}$ в
стационаре, а $n_{\rm out}$ поглощений уносят $n_{\rm out}\msink$, --- так
что
\begin{equation}
  E \;\simeq\; N\sqrt{\msink} \;+\; n_{\rm out}\,\msink .
  \label{eq:num_budget}
\end{equation}
Второе слагаемое не зависит от $N$: ворота не дорожают с ростом прогона,
дорожает только заполнение ящика. Уравнение~\eqref{eq:num_budget} выведено
при $\lambda = 0$ и переоценивает стоимость при $\lambda > 0$, где более
крутое ядро гонит массу к стоку быстрее в расчёте на событие.

У замкнутой системы стока нет, поэтому взведённое значение по умолчанию там
молча разоружается; исключение бросается только на явно переданное
значение. Различие между <<это набрал пользователь>> и <<это пришло из
сигнатуры>> несёт подкласс \code{int}, так что во всех прочих местах
значение ведёт себя как обычное целое.

\section{Измерительный слой}
\label{sec:num_measure}

Прогон производит массивы; превращение их в показатель степени --- отдельная
стадия со своими режимами отказа.

\paragraph{Оценщик.} Замкнутая система: про\-ин\-те\-гри\-ро\-ван\-ная по
возрасту суперпозиция $F(m) = \sum_t (dN/dm)\,\Delta t$. Открытая система:
стационарное состояние \emph{и есть} мгновенный спектр, поэтому оценщиком
служит последний снапшот, а $\Delta t$ в него не входит.

\paragraph{Локальный наклон и плато.} Честный способ смотреть на измеренный
спектр --- это
\begin{equation}
  \Gamma(m) \;=\; \frac{d\log F}{d\log m},
\end{equation}
вычисленный методом наименьших квадратов на скользящем окне. Настоящий
инерционный интервал --- это \emph{плато} в $\Gamma$; загрязнённые концы
проявляются как места, где $\Gamma$ загибается. Сообщаемый индекс есть
среднее $\Gamma$ по самому длинному непрерывному участку, на котором оно
остаётся плоским в пределах допуска, вместе с его разбросом и шириной в
декадах. Плато уже примерно одной декады не является степенным законом,
сколь бы узкой ни была его погрешность; ширина печатается рядом с индексом
именно поэтому. Если ни один участок не проходит, ответ --- \code{nan}, и
это правильный ответ, а не отказ.

\paragraph{Защитная полоса (guard band).} Независимо от этого: отрезать по
$0{,}5$~декады с каждого конца интервала между двумя характерными массами
задачи --- $[\minj,\msink]$ для открытой системы, $[m_i, m_0^{\max}]$ для
замкнутой --- и подогнать там единственный степенной закон. Эта полоса
выбирается \emph{до} того, как посмотрели на данные. Две оценки обязаны
согласоваться; если нет, каскад не развился на том диапазоне, который был
предположен.

Подгонка по всему массиву не даёт ни того, ни другого. На одном прогоне
замкнутой коагуляции при $\lambda=0$ и теоретическом $\alpha=-1$: весь
массив $-1{,}71$, защитная полоса $-1{,}09$, плато $-1{,}16 \pm 0{,}15$ на
$1{,}3$ декады.

\paragraph{Компенсированное представление.} На панелях рисуется
\begin{equation}
  m^2 \frac{dN}{dm} \;=\; \text{масса на логарифмический интервал массы},
\end{equation}
чей наклон равен $\alpha+2$. Вся подгонка делается по сырому $dN/dm$;
компенсация применяется только при отрисовке. Модельная линия несёт индекс
плато и амплитуду, полученную как свободный член наименьших квадратов при
\emph{фиксированном} наклоне, так что никакая вторая подгонка не может
тихо сдвинуть показатель, а теоретическая линия привязывается тем же
способом --- тогда две кривые различаются только наклоном.

\paragraph{Предсказание.} При степени однородности ядра $\lambda$
\begin{equation}
  \beta =
    \begin{cases}
      \lambda & \text{замкнутая: пакет } n \propto m_0^{-2},\\[2pt]
      \dfrac{1+\lambda}{2} & \text{открытая: фон } n \propto m_0^{\alpha},
    \end{cases}
  \qquad
  \alpha = -(1+\beta),
  \quad
  b = \frac{1}{1-\beta}.
\end{equation}
Коагуляция и фрагментация разделяют $\beta$, $b$ и $\alpha$ в точности;
переворачивается только знак дрейфа, поэтому закон читается в $\tau$ при
движении вверх и в $\tau_\ast - \tau$ при движении вниз.

\begin{center}
\footnotesize
\begin{tabular}{@{}@{}l@{\ \ }l@{\ \ }c@{\ \ }c@{\ \ }c@{}}
\toprule
ядро & $K(m_1,m_2)$ & $\lambda$ &
  замк.\ $\alpha$, $b$ &
  откр.\ $\alpha$, $b$ \\
\midrule
постоянное     & $1$                                & $0$   & $-1$,\ \ $1$    & $-3/2$,\ \ $2$ \\
сумма          & $m_1^{1/3}+m_2^{1/3}$              & $1/3$ & $-4/3$,\ \ $3/2$& $-5/3$,\ \ $3$ \\
геометрич.\    & $(m_1^{1/3}+m_2^{1/3})^2$          & $2/3$ & $-5/3$,\ \ $3$  & $-11/6$,\ \ $6$ \\
аддитивное     & $m_1+m_2$                          & $1$   & $-2$,\ \ $\infty$ & $-2$,\ \ $\infty$ \\
произвед.\     & $m_1 m_2$                          & $2$   & \multicolumn{2}{c}{выше гелеобразования} \\
\bottomrule
\end{tabular}
\end{center}

\section{Валидация}
\label{sec:num_validation}

Постоянное ядро, $\lambda = 0$. \emph{плато} --- автоматический индекс,
\emph{полоса} --- априорная защитная полоса, \emph{дек} --- ширина плато в
декадах.

\begin{center}
\small
\begin{tabular}{@{}l r r r r r r@{}}
\toprule
случай & $b$ & $b_{\rm теор}$ & плато & $\pm$ & дек & полоса \\
\midrule
1 замкнутая коагуляция     & $0{,}990$ & $1$ & $-1{,}077$ & $0{,}126$ & $2{,}70$ & $-1{,}035$ \\
2 замкнутая фрагментация   & $1{,}001$ & $1$ & $-0{,}978$ & $0{,}141$ & $3{,}00$ & $-1{,}056$ \\
3 открытая коагуляция      & $2{,}001$ & $2$ & $-1{,}475$ & $0{,}095$ & $3{,}60$ & $-1{,}498$ \\
4 открытая фрагм., $f=0$   & ---       & $2$ & $-2{,}047$ & $0{,}183$ & $0{,}80$ & $-2{,}354$ \\
4 открытая фрагм., $f=0{,}3$ & ---     & $2$ & $-1{,}550$ & $0{,}091$ & $3{,}40$ & $-1{,}517$ \\
\bottomrule
\end{tabular}
\end{center}

Случай~1 дополнительно воспроизводит точное решение Смолуховского для
постоянного ядра, $m_0(t) = m_i(1+n_0K_1t/2)$, с точностью
$\max|{\rm отношение}-1| = 7{,}8\times10^{-14}$ --- прямая проверка
уравнения~\eqref{eq:num_clock}. Дрейф массы тождественно нулевой во всех
четырёх случаях.

Две строки следует читать как диагноз, а не как измерение. Случай~4 при
$f=0$ возвращает $-2{,}05$ на $0{,}80$ декады: индекс залип на $-2$ в
точности так, как предсказывает разд.~\ref{sec:trick8}, а узкое плато само
объявляет, что это не измерение. И показатель роста $b$ для случая~4
отсутствует потому, что при таком бюджете более $99\%$ живых частиц всё ещё
происходят от начального условия и наследуют $t_{\rm born}=0$; ось возраста
тогда вырождена --- возраст равен астрономическому времени сразу у всех, ---
и изохроны деградируют в <<спектр как функция времени>>. Тот же факт
проявляется как $M_{\rm out}/M_{\rm in} > 1$: сток сливает начальный
резервуар, чья масса в $M_{\rm in}$ никогда не входила.

\section{Практические режимы отказа}
\label{sec:num_pitfalls}

\paragraph{Масса вне сетки --- жёсткое падение.}
Мажоранта строится по \emph{верхним углам} бинов, а $b(m)$ зажимает всё,
что выше верхней границы, в последний бин, чей угол тогда меньше самой
частицы. Следовательно $K(m,m) > \fmaj_{B-1,B-1}$, и первая же пара бросает
\code{Majorant violated}. При постоянном ядре этого не может случиться в
принципе, поскольку $K\equiv1$ плоское; при любом $\lambda>0$ это
происходит немедленно. Проверяйте, что сетка охватывает $\minj$ и $\msink$,
до запуска, а не после.

\paragraph{Неограниченный бюджет событий безопасен ровно в одном случае.}
Снимать ограничение на число событий можно только тогда, когда прогон
способно завершить что-то другое:

\begin{center}
\small
\begin{tabular}{@{}l l p{7.4cm}@{}}
\toprule
случай & без предела & что останавливает --- или не останавливает \\
\midrule
замкнутая коагуляция   & безопасно & \code{stop\_max\_mass}; к тому же $N$ монотонно
                                     убывает, так что \code{live}~$<2$ --- гарантированный
                                     терминатор \\
замкнутая фрагментация & опасно    & $N$ растёт; сторож зависания срабатывает лишь
                                     тогда, когда \emph{всё} окажется ниже пассивного
                                     пола, то есть при $N_0 m_i/m_{\rm floor}$ частицах \\
открытая коагуляция    & зависнет  & $N$ сидит на полке, поэтому никакое другое
                                     условие сработать не может \\
открытая фрагментация  & опасно    & $N$ растёт неограниченно \\
\bottomrule
\end{tabular}
\end{center}

В открытых случаях правильный тормоз --- \code{stop\_sink\_events}
(разд.~\ref{sec:num_sink}); дополнительно оставленный ресурсный потолок не
стоит ничего и ограничивает последствия неверно настроенной инжекции.
Записанная причина остановки затем говорит, какой из двух сработал.

\paragraph{Смена ядра обесценивает четыре константы.}
Однородность $\lambda$ и все предсказанные показатели подхватываются
автоматически, а искатель плато, защитная полоса и компенсация вообще не
знают про ядро. Четыре вещи --- знают:

\begin{enumerate}
  \item \textbf{Темп инжекции.} Баланс числа в стационаре есть
        $q = \langle K\rangle N^2/(2V)$. Запись $q = N^2/2$ --- это та же
        формула при $\langle K\rangle = 1$. Поскольку $N$ набирается на
        масштабе инжекции при $\alpha<-1$, большинство пар суть
        $(\minj,\minj)$, так что $\langle K \rangle \sim K(\minj,\minj)$, а
        более тяжёлые партнёры поднимают эту величину примерно вдвое.
  \item \textbf{Время столкновения.} $t_c = 1/(Kn)$, а не $1/n$.
  \item \textbf{Биновка по возрасту.} Один возрастной бин шириной
        $\Delta_\tau$ декад покрывает $b\,\Delta_\tau$ декад \emph{массы}.
        Фиксируйте разрешение по массе и выводите шаг по возрасту из него,
        $\Delta_\tau = 0{,}30/b$; иначе при $b=6$ шаг, настроенный под
        $b=2$, перепрыгивает почти декаду массы за бин, и двухдекадное окно
        подгонки ловит два бина.
  \item \textbf{Точные решения.} Замкнутая форма Смолуховского и проверка
        <<$m_0(t)$ обязано быть линейным>> --- свойства исключительно
        постоянного ядра.
\end{enumerate}

\paragraph{У отпечатка $-2$ три разные причины.}
Индекс $-2$ возвращают: (i) настоящее ядро с $\lambda = 1$, где это
правильный ответ; (ii) нелокальное правило разрушения из
разд.~\ref{sec:trick8}, где это граничный эффект; (iii) часы, продвигаемые
на событие, а не на попытку, разд.~\ref{sec:trick5}, где это численный
артефакт. Все три гонят $\beta \to 1$ и на бумаге выглядят одинаково. Само
по себе значение $-2$ поэтому не говорит ни о чём. Различающая проверка ---
меняется ли индекс при смене ядра: если не меняется, это не универсальность.

\paragraph{Плоская полка --- не инерционный интервал.}
В замкнутой фрагментации пассивный пол замораживает всё, что ниже него: эти
частицы больше никогда не дробятся и накапливаются в плоскую полку, к
каскаду отношения не имеющую. Оставленная в массиве, эта полка --- самый
длинный плоский участок из всех, $4{,}3$ декады при $\alpha \simeq 0$, --- и
искатель плато послушно сообщает именно о ней. Ограничивайте поиск областью
$m \ge m_{\rm floor}$.

% ---------------------------------------- удалить при вставке в main.tex
\end{document}
